\documentclass[11pt,a4paper,oneside]{book}
\usepackage[T1]{fontenc}
\usepackage[utf8]{inputenc}
\usepackage[indonesian]{babel}
\usepackage{lmodern}
\usepackage{amsmath,amsthm,amssymb,amscd}
\usepackage{graphicx}
\usepackage[a4paper,margin=22mm,headheight=15pt]{geometry}
\usepackage{microtype}
\usepackage[hidelinks,pdfusetitle]{hyperref}
\usepackage{bookmark}
% language package supplied by unit wrapper

% UTF-8 input supplied by unit wrapper

%\usepackage[utf8]{inputenc}

\usepackage{ifthen}

%\usepackage{fancyhdr}

%Aus Grundvorspann

%Packages

\usepackage{amscd}

\usepackage{amssymb}

\usepackage{epsfig} %Graphik

\usepackage{verbatim,moreverb} %

\usepackage{enumitem}

\usepackage{eurosym}

\usepackage{epigraph}

\setcounter{MaxMatrixCols}{20}

%Mathematische Einzelbefehle

%[[Projekt:Semantische Vorlagen/Mathematischer Modus/Latex]]

\newcommand{\mathdisp}[2]{\[ #1 \, #2 \]}

\newcommand{\mathdisplayteile}[3]{\[ #1 \] \[ #2 \, #3 \]}

\newcommand{\mathind}[3]{$ #1$, $#2 $#3}

\newcommand{\mathbed}[8]{$ #1 ${\ifthenelse{\equal{#2}{}}{, }{ #2 }}$ #3 ${\ifthenelse{\equal{#5}{}}{}{{\ifthenelse{\equal{#4}{}}{, }{ #4 }}}$ #5 $}{\ifthenelse{\equal{#7}{}}{}{{\ifthenelse{\equal{#6}{}}{, }{ #6 }}}$ #7 $}#8}

\newcommand{\mathbeddisp}[8]{\[ #1 {\ifthenelse{\equal{#2}{}}{,\, }{ \text{ #2 } }} #3 {\ifthenelse{\equal{#5}{}}{}{{\ifthenelse{\equal{#4}{}}{,\, }{ \text{ #4 } }}} #5 }{\ifthenelse{\equal{#7}{}}{}{{\ifthenelse{\equal{#6}{}}{,\, }{ \text{ #6 } }}} #7 }#8 \]}

\newcommand{\mathbedtermdisp}[8]{\[ #1\] #2 $ #3 $#4 $ #5 $#6$ #7 $#8}

\newcommand{\mathkon}[4]{$ #1 $ #2 $ #3 $#4}

\newcommand{\mathkor}[5]{{\ifthenelse{\equal{#1}{}}{}{#1 }}$ #2 $ #3 $ #4 $#5}

\newcommand{\mathkork}[5]{{\ifthenelse{\equal{#1}{}}{}{#1 }}$ #2 $ #3 $ #4 $#5}

\newcommand{\mathlist}[5]{$ #1 $\ifthenelse{\equal{#2}{}}{,}{ #2} $ #3 $\ifthenelse{\equal{#4}{}}{,}{ #4} $ #5 $}

\newcommand{\mathlistdisp}[5]{\[  #1  \ifthenelse{\equal{#2}{}}{,}{ \text{ #2 }}  #3  \ifthenelse{\equal{#4}{}}{,}{\text{ #4 } }  #5 \] }

\newcommand{\alignier}[2]{\begin{align*} #1 \, #2 \end{align*}}

\newcommand{\mathlistvierdisp}[7]{\[ #1,\, #3, \, #5 \text{ #6 } #7  \]}

%[[Projekt:Semantische Vorlagen/Mathematische Vergleichsvorlagen/Latex]]

\newcommand{\mavergleichskette}[4]{$ #1 #2 #3#4$}

\newcommand{\mavergleichskettedisp}[4]{\[  #1 #2 #3 #4 \]}

\newcommand{\mavergleichskettedisplang}[4]{\[  #1 #2 #3 #4 \]}

\newcommand{\mavergleichskettealign}[4]{ \begin{eqnarray*}  #1 #2 #3 #4 \end{eqnarray*}  }

\newcommand{\vergleichskette}[9]{ #1 \, #2 \, #3 \ifthenelse
{\equal {#4}{ }} {} {\, #4 \, #5 }  \ifthenelse {\equal {#6}{ } }{}
{\, #6 \, #7 } \ifthenelse {\equal {#8}{ }}{} {\, #8 \, #9 }  }

\newcommand{\vergleichskettelang}[9]{ #1 \, #2 \, #3 \ifthenelse
{\equal {#4}{ }} {} {\, #4 \, #5 }  \ifthenelse {\equal {#6}{ } }{}
{\, #6 \, #7 } \ifthenelse {\equal {#8}{ }}{} {\, #8 \, #9 }  }

\newcommand{\vergleichskettefortsetzung}[8]{ \, #1 \, #2 \ifthenelse
{\equal {#3}{ }} {} {\, #3 \, #4 }  \ifthenelse {\equal {#5}{ } }{}
{\, #5 \, #6 } \ifthenelse {\equal {#7}{ }}{} {\, #7 \, #8 } }

\newcommand{\vergleichskettealign}[9]{ #1 & #2 & #3 \ifthenelse
{\equal {#4}{ }} {} {\cr & #4 & #5 }  \ifthenelse {\equal {#6}{ }
}{} {\cr & #6 & #7 } \ifthenelse {\equal {#8}{ }}{} {\cr & #8 & #9 }
}

\newcommand{\vergleichskettefortsetzungalign}[8]{ \cr & #1 & #2 \ifthenelse
{\equal {#3}{ }} {} {\cr & #3 & #4 }  \ifthenelse {\equal {#5}{ }
}{} {\cr & #5 & #6 } \ifthenelse {\equal {#7}{ }}{} {\cr & #7 & #8 }
}

\newcommand{\mavergleichskettealigndrucklinks}[4]{ \begin{eqnarray*}  #1 #2 #3 #4 \end{eqnarray*}  }

\newcommand{\vergleichskettealigndrucklinks}[9]{  &  & #1 \cr & #2 & #3 \ifthenelse
{\equal {#4}{ }} {} {\cr & #4 & #5 }  \ifthenelse {\equal {#6}{ }
}{} {\cr & #6 & #7 } \ifthenelse {\equal {#8}{ }}{} {\cr & #8 & #9 }
}

\newcommand{\mavergleichskettealignhandlinks}[4]{ \begin{eqnarray*}  #1 #2 #3 #4 \end{eqnarray*}  }

\newcommand{\vergleichskettealignhandlinks}[9]{ #1 & #2 & #3 \ifthenelse
{\equal {#4}{ }} {} {\cr & #4 & #5 }  \ifthenelse {\equal {#6}{ }
}{} {\cr & #6 & #7 } \ifthenelse {\equal {#8}{ }}{} {\cr & #8 & #9 }
}

\newcommand{\mavergleichskettedisphandlinks}[4]{ \[ #1 #2 #3 #4 \] }

\newcommand{\vergleichskettedisphandlinks}[9]{  #1 \, #2 \, #3 \ifthenelse
{\equal {#4}{ }} {} {\, #4 \, #5 }  \ifthenelse {\equal {#6}{ } }{}
{\, #6 \, #7 } \ifthenelse {\equal {#8}{ }}{} {\, #8 \, #9 }  }

\newcommand{\zeilemitteil}[2]{ \ifthenelse {\equal {#2} {} } {#1} { #1 \cr & & #2}   }

\newcommand{\mavergleichskettek}[4]{$ #1 #2 #3#4$}

\newcommand{\vergleichskettek}[9]{ #1 \, #2 \, #3 \ifthenelse
{\equal {#4}{ }} {} {\, #4 \, #5 }  \ifthenelse {\equal {#6}{ } }{}
{\, #6 \, #7 } \ifthenelse {\equal {#8}{ }}{} {\, #8 \, #9 }  }

\newcommand{\vergleichskettefortsetzungk}[8]{ \, #1 \, #2 \ifthenelse
{\equal {#3}{ }} {} {\, #3 \, #4 }  \ifthenelse {\equal {#5}{ } }{}
{\, #5 \, #6 } \ifthenelse {\equal {#7}{ }}{} {\, #7 \, #8 } }

%[[Projekt:Semantische Vorlagen/Mathematische Strukturvorlagen/Latex]]

\newcommand{\N}   {{\mathbb N}}
\newcommand{\Z}   {{\mathbb Z}}
\newcommand{\Q}   {{\mathbb Q}}
\newcommand{\R}   {{\mathbb R}}
\newcommand{\C}   {{\mathbb C}}
\newcommand{\Complex}  {{\mathbb C}}

\newcommand{\defeq}{:=}

\newcommand{\defeqr}{=:}

\newcommand{\maabb}[4]{${\ifthenelse {\equal {#1}{} } {} {#1 \colon }}\,  #2 \rightarrow #3 #4 $}

\newcommand{\maabbele}[6]{$ {\ifthenelse {\equal {#1}{} } {} {#1 \colon }}\,   #2 \rightarrow  #3 , \, #4 \mapsto #5 #6 $}

\newcommand{\maabbdisp}[4]{\[ {\ifthenelse {\equal {#1}{} } {} {#1 \colon }}\,  #2 \longrightarrow #3 #4 \]}

\newcommand{\maabbeledisp}[6]{\[ {\ifthenelse {\equal {#1}{} } {} {#1 \colon }}\,   #2 \longrightarrow  #3 , \, #4 \longmapsto #5 #6 \]}

%Variante bei ifequal-Problem

\newcommand{\maabbnamedisp}[4]{\[ #1 \colon \, #2 \longrightarrow #3 #4 \]}

\newcommand{\maabbeledispvar}[6]{\[ {\ifthenelse {\equal {#1}{} } {} {#1 \colon }}\,   #2 \longrightarrow  #3 , \, #4 \longmapsto #5 #6  \]}

\newcommand{\maabbelementzeiledisplay}[6]{\[ {\ifthenelse {\equal {#1}{} } {} {#1 \colon }}\, #2 \longrightarrow #3 , \]  	\[ #4 \longmapsto #5 #6 \]}

\newcommand{\maabbelementdoppelzeiledisplay}[6]{\[ {\ifthenelse {\equal {#1}{} } {} {#1 \colon }}\, #2 \longrightarrow #3 , \] \[ #4 \longmapsto \] \[  #5 #6 \]}

\newcommand{ \arccot  } {\operatorname{arccot } }

\newcommand{\betrag}[1]{\left| #1 \right|}

\newcommand{\strukturgarbe}{ {\mathcal O} }

\newcommand{\schnittring}[1]{\Gamma (#1 , {\mathcal O})}

\newcommand{\tensor}{\otimes}

%[[Projekt:Semantische Vorlagen/Textgestaltung/Latex]]

\newcommand{\einrueckung}[1]{\par \hspace{1cm}#1 \par}

\newcommand{\zusatz}[2]{(#1)#2}

\newcommand{\zusatzklammer}[3]{(#1#2)#3}

\newcommand{\zusatzgs}[3]{\-- #1 \-- #3}

\newcommand{\zusatzfussnote}[3]{#3\footnote{#1#2}}

%Feinabstimmung

\newcommand{\latexdruckeinskurz}{\!}

\newcommand{\latexdruckzweikurz}{\! \!}

\newcommand{\latexdruckdreikurz}{\! \! \!}

%[[Projekt:Semantische Vorlagen/Satzumgebungen/Latex]]

\newtheorem{fakt}{Fakta}[section] %so wird abschnittsweise %durchnummeriert. Bei [] wird % einfach nummeriert.

\newtheorem{Fakt}{fakt}[Fakt]

\newtheorem{Proposition}[fakt]{Proposisi}

\newtheorem{Korollar}[fakt]{Korolari}

\newtheorem{Lemma}[fakt]{Lema}

\newtheorem{Satz}[fakt]{Teorema}

\newtheorem{Theorem}[fakt]{Teorema}

\theoremstyle{definition}

\newtheorem{Axiom}[fakt]{Aksioma}

\newtheorem{aufgabe}[fakt]{Soal}

\newtheorem{Aufgabe}[fakt]{Soal}

\newtheorem{klausuraufgabe}{Soal ujian}

\newtheorem{Klausuraufgabe}{Soal ujian}

\newtheorem{beispiel}[fakt]{Contoh}

\newtheorem{Beispiel}[fakt]{Contoh}

\newtheorem{Konstruktion}[fakt]{Konstruksi}

\newtheorem{konstruktion}[fakt]{Konstruksi}

\newtheorem{Verfahren}[fakt]{Verfahren}

\newtheorem{verfahren}[fakt]{Verfahren}

\newtheorem{bemerkung}[fakt]{Catatan}

\newtheorem{Bemerkung}[fakt]{Catatan}

\newtheorem{definition}[fakt]{Definisi}

\newtheorem{Definition}[fakt]{Definisi}

\newtheorem{notation}[fakt]{Notation}

\newtheorem{Notation}[fakt]{Notation}

\newtheorem{Frage}[fakt]{Frage}

\newtheorem{frage}[fakt]{Frage}

\newtheorem{problem}[fakt]{Problem}

\newtheorem{Problem}[fakt]{Problem}

\newtheorem{situation}[fakt]{Situation}

\newtheorem{Situation}[fakt]{Situation}

%Einlesegestaltung

%[[Projekt:Semantische Vorlagen/Umgebungsdesign/Latex]]

%Aufgabengestaltung

\newcommand{\inputaufgabe}[4]{\aufgabevorskip \begin{aufgabe} \punkte {#1} \aufgabepunktskip #2 #3 #4\end{aufgabe} \aufgabenachskip}

\newcommand{\inputaufgabegibtloesung}[4]{\aufgabevorskip \begin{aufgabe}\hspace{-1.8 mm}* \punkte {#1} \aufgabepunktskip #2 #3 #4\end{aufgabe} \aufgabenachskip}

\newcommand{\inputaufgabepunktenummervonhand}[3]{\aufgabevorskip {\bf Aufgabe #1} (#2 Punkte) \aufgabepunktskip #3  \aufgabenachskip}

\newcommand{\inputaufgabeloesung}[2]{ \begin{aufgabe} #1 \end{aufgabe} \aufgabenachskip

\bigskip

Lösung

\bigskip #2

}

\newcommand{\inputaufgabeklausurloesung}[3]{\par \newpage \begin{aufgabe} {\ifthenelse {\equal {#1}{}}{} {\ifthenelse {\equal {#1}{1}} {(1 Punkt)} {(#1 Punkte)}}} \par \bigskip #2 \end{aufgabe} \underline{L\"osung:} \par \bigskip #3 }

\newcommand{\inputaufgabeloesungvar}[2]{\aufgabevorskip Aufgabe: #1 \aufgabenachskip
Lösung: #2}

\newcommand{\inputaufgabepunkteloesung}[3]{\aufgabevorskip \begin{aufgabe} {\ifthenelse {\equal {#1}{}}{} {\ifthenelse {\equal {#1}{1}} {(1 Punkt)} {(#1 Punkte)}}}  \aufgabepunktskip #2 \end{aufgabe} \aufgabenachskip \loesung #3}

\newcommand{\loesung}[1]{\underline{Lösung:} #1}

\newcommand{\punkte}[1]{\ifthenelse {\equal {#1}{}}{} {\ifthenelse {\equal {#1}{1}} {(1 poin)} {(#1 poin)}}}

\newcommand{\inputexercise}[4]{\aufgabevorskip \begin{exercise} \points {#1} \aufgabepunktskip #2 #3 #4\end{exercise} \aufgabenachskip}

\newcommand{\points}[1]{\ifthenelse {\equal {#1}{}}{} {\ifthenelse {\equal {#1}{1}} {(1 point)} {(#1 points)}}}

%Beispielgestaltung

\newcommand{\inputbeispiel}[2]
{\beispielvorskip
\begin{beispiel}
\beispielbenennung{#1}
\beispielbenennungnachskip #2
\end{beispiel}
\beispielnachskip}

\newcommand{\beispielbenennung}[1]{ \ifthenelse {\equal {#1}{} \OR \equal {#1}{   }  }{}{(#1)}     }

%Bemerkunggestaltung

\newcommand{\inputbemerkung}[2]
{\bemerkungvorskip
\begin{bemerkung}
\bemerkungbenennung{#1}
\bemerkungbenennungnachskip #2
\end{bemerkung}
\bemerkungnachskip}

\newcommand{\bemerkungbenennung}[1]{  \ifthenelse {\equal {#1}{} \OR \equal {#1}{   }  }{}{(#1)}}

\newcommand{\inputverfahren}[2]
{\bemerkungvorskip
\begin{verfahren}
\bemerkungbenennung{#1}
\bemerkungbenennungnachskip #2
\end{verfahren}
\bemerkungnachskip}

\newcommand{\inputkonstruktion}[2]
{\bemerkungvorskip
\begin{konstruktion}
\bemerkungbenennung{#1}
\bemerkungbenennungnachskip #2
\end{konstruktion}
\bemerkungnachskip}

\newcommand{\inputfrage}[2]
{\bemerkungvorskip
\begin{frage}
\problembenennung{#1}
\bemerkungbenennungnachskip #2
\end{frage}
\bemerkungnachskip}

\newcommand{\fragebenennung}[1]{  \ifthenelse {\equal {#1}{} \OR \equal {#1}{   }  }{}{(#1)}}

\newcommand{\inputproblem}[2]
{\bemerkungvorskip
\begin{problem}
\problembenennung{#1}
\bemerkungbenennungnachskip #2
\end{problem}
\bemerkungnachskip}

\newcommand{\problembenennung}[1]{  \ifthenelse {\equal {#1}{} \OR \equal {#1}{   }  }{}{(#1)}}

%Situationsgestaltung

\newcommand{\inputsituation}[2]
{\begin{situation}
\situationbenennung{#1} #2
\end{situation}
}

\newcommand{\situationbenennung}[1]{\ifthenelse {\equal {#1}{} \OR \equal {#1}{  }  }{}{(#1)}}

%Definitiongestaltung

\newcommand{\inputdefinition}[2]
{\definitionvorskip
\begin{definition}
\definitionbenennung{#1}
\definitionbenennungnachskip #2
\end{definition}
\definitionnachskip}

\newcommand{\inputaxiom}[2]
{\definitionvorskip
\begin{Axiom}
\definitionbenennung{#1}
\definitionbenennungnachskip #2
\end{Axiom}
\definitionnachskip}

\newcommand{\definitionbenennung}[1]{\ifthenelse {\equal {#1}{} \OR \equal {#1}{  }  }{}{(#1)}}

\newcommand{\inputnotation}[2]
{\definitionvorskip \begin{notation} \definitionbenennung{#1} \definitionbenennungnachskip #2 \end{notation} \definitionnachskip}

%Fakt- und Beweisgestaltung

\renewcommand{\proofname}{\hspace{0cm}{\it Bukti}}

\newcommand{\beweis}[1]{\begin{proof} #1 \end{proof}}

\newcommand{\inputfakt}[4]
{\faktvorskip
\begin{#2}

%\label{#1}   (Absetzen, sonst  kann das folgende dahinter rutschen)

\faktbenennung{#3}
\faktbenennungnachskip #4
\end{#2}
\faktnachskip
}

\newcommand{\inputfaktbeweis}[5]
{\faktvorskip
\begin{#2}

%\label{#1}   (Absetzen, sonst  kann das folgende dahinter rutschen)

\faktbenennung{#3}
\faktbenennungnachskip #4
\end{#2}
\faktnachskip
\beweis{#5}}

\newcommand{\inputfaktbeweisnichtvorgefuehrt}[5]{\inputfaktbeweis {#1} {#2} {#3} {#4} {Dieser Beweis wurde in der Vorlesung nicht vorgef\"uhrt.} }

\newcommand{\inputfaktbeweistrivial}[4]{\inputfaktbeweis {#1} {#2} {#3} {#4} {Das ist trivial.} }

\newcommand{\inputfaktuebergangbeweis}[6]
{\faktvorskip
\begin{#2}

%\label{#1} (Absetzen, sonst kann das
% folgende dahinter rutschen)

\faktbenennung{#3}
\faktbenennungnachskip #4
\end{#2}
\faktnachskip #5 \faktnachskip
\beweis{#6}}

\newcommand{\inputbeweis}[1]{\beweis{#1} }

\newcommand{\faktbenennung}[1]{\ifthenelse {\equal {#1}{} \OR \equal {#1}{  }  }{}{(#1)}}

%Gestaltung der Faktstruktur

\newcommand{\faktsituation}[1]{\faktsituationskip #1}

\newcommand{\faktvoraussetzung}[1]{\faktvoraussetzungskip #1}

\newcommand{\faktvoraussetzungleer}[1]{ \hspace{-0,15cm} }

\newcommand{\faktvoraussetzungpos}[1]{\faktvoraussetzungskip #1}

\newcommand{\faktuebergang}[1]{\faktuebergangskip #1}

\newcommand{\faktuebergangleer}[1]{\faktuebergangskip \hspace{-0,15cm} }

\newcommand{\faktfolgerung}[1]{\faktfolgerungskip #1}

\newcommand{\faktzusatz}[1]{\faktzusatzskip #1}

%[[Projekt:Semantische Vorlagen/Beweisgestaltung/Latex]]

\newcommand{\teilbeweis}[5]{#1#2#3#4#5}

\newcommand{\fallunterscheidung}[5]{#1 #2\ifthenelse {\equal {#3}{}} {} { #3}\ifthenelse {\equal {#4}{}} {} { #4}\ifthenelse {\equal {#5}{}} {} { #5}}

\newcommand{\fallunterscheidungzwei}[2]{#1 #2}

\newcommand{\fallunterscheidungdrei}[3]{#1 #2 #3 }

\newcommand{\fallunterscheidungvier}[4]{#1 #2 #3 #4 }

\newcommand{\fallunterscheidungfuenf}[5]{#1 #2 #3 #4 #5 }

%[[Projekt:Semantische Vorlagen/Stichwortumsetzung/Latex]]

\newcommand{\betonung}[2]{\emph{#1}#2}

\newcommand{\stichwort}[2]{\emph{#1}#2}

\newcommand{\stichwortpraemath}[3]{$#1$-\emph{#2}#3}

\newcommand{\definitionswort}[2]{\emph{#1}#2}

\newcommand{\definitionswortpraemath}[3]{$#1$-\emph{#2}#3}

\newcommand{\definitionswortteil}[2]{\emph{#1}#2}

\newcommand{\definitionswortenp}[2]{\emph{#1}#2}

\newcommand{\definitionsverweis}[3]{#1#3}

\newcommand{\definitionsverweisanfuehrung}[3]{\anfuehrung {#1} {#3} }

%Anführungszeichen

\newcommand{\anfuehrung}[2]{"`{#1}"'#2}

\newcommand{\anfuehrungenglisch}[2]{``{#1}''#2}

%[[Projekt:Semantische Auflistungsvorlagen/Latex]]

%Aufzählungen

\newcommand{\aufzaehlungeins}[1]{\begin{enumerate} \item  #1 \end{enumerate}}

\newcommand{\aufzaehlungzwei}[2]{\begin{enumerate} \item  #1 \item  #2 \end{enumerate}}

\newcommand{\aufzaehlungdrei}[3]{\begin{enumerate} \item  #1 \item  #2 \item  #3  \end{enumerate}}

\newcommand{\aufzaehlungvier}[4]{\begin{enumerate} \item  #1 \item  #2 \item  #3 \item  #4 \end{enumerate}}

\newcommand{\aufzaehlungfuenf}[5]{\begin{enumerate} \item  #1 \item  #2 \item  #3 \item  #4 \item  #5  \end{enumerate}}

\newcommand{\aufzaehlungsechs}[6]{\begin{enumerate} \item  #1 \item  #2 \item  #3 \item  #4 \item  #5 \item  #6 \end{enumerate}}

\newcommand{\aufzaehlungsieben}[7]{\begin{enumerate} \item  #1 \item  #2 \item  #3 \item  #4 \item  #5 \item #6 \item  #7 \end{enumerate}}

\newcommand{\aufzaehlungacht}[8]{\begin{enumerate} \item  #1 \item  #2 \item  #3 \item  #4 \item  #5 \item #6 \item  #7 \item #8 \end{enumerate}}

\newcommand{\aufzaehlungneun}[9]{\begin{enumerate} \item  #1 \item  #2 \item  #3 \item  #4 \item  #5 \item #6 \item  #7  \item  #8  \item  #9 \end{enumerate}}

\newcommand{\itemfuenf}[5]{\item  #1 \item  #2 \item  #3 \item  #4 \item  #5}

\newcommand{\itemsechs}[6]{\item  #1 \item  #2 \item  #3 \item  #4 \item  #5 \item  #6}

\newcommand{\itemsieben}[7]{\item  #1 \item  #2 \item  #3 \item  #4 \item  #5 \item  #6 \item #7}

\newcommand{\aufzaehlungeinsabc}[1]{\begin{enumerate}[label=(\alph*)]  #1 \end{enumerate}}

\newcommand{\aufzaehlungzweiabc}[2]{\begin{enumerate}[label=(\alph*)] \item  #1 \item  #2 \end{enumerate}}

\newcommand{\aufzaehlungdreiabc}[3]{\begin{enumerate}[label=(\alph*)] \item  #1 \item  #2 \item  #3  \end{enumerate}}

\newcommand{\aufzaehlungvierabc}[4]{\begin{enumerate}[label=(\alph*)] \item  #1 \item  #2 \item  #3 \item  #4 \end{enumerate}}

\newcommand{\aufzaehlungfuenfabc}[5]{\begin{enumerate}[label=(\alph*)] \item  #1 \item  #2 \item  #3 \item  #4 \item  #5  \end{enumerate}}

\newcommand{\aufzaehlungsechsabc}[6]{\begin{enumerate}[label=(\alph*)] \item  #1 \item  #2 \item  #3 \item  #4 \item  #5 \item  #6 \end{enumerate}}

\newcommand{\aufzaehlungsiebenabc}[7]{\begin{enumerate}[label=(\alph*)] \item  #1 \item  #2 \item  #3 \item  #4 \item  #5 \item #6 \item  #7 \end{enumerate}}

\newcommand{\aufzaehlungachtabc}[8]{\begin{enumerate}[label=(\alph*)] \item  #1 \item  #2 \item  #3 \item  #4 \item  #5 \item #6 \item  #7 \item #8 \end{enumerate}}

\newcommand{\aufzaehlungneunabc}[9]{\begin{enumerate}[label=(\alph*)] \item  #1 \item  #2 \item  #3 \item  #4 \item  #5 \item #6 \item  #7  \item  #8  \item  #9 \end{enumerate}}

\newcommand{\aufzaehlungzweireihe}[2]{\begin{enumerate} #1 #2 \end{enumerate}}

\newcommand{\aufzaehlungdreibeweis}[3]{(1) #1 \par (2) #2 \par (3) #3 }

\newcommand{\aufzaehlungvierbeweis}[4]{(1) #1 \par (2) #2 \par (3) #3  \par (4) #4}

\newcommand{\aufzaehlungfuenfbeweis}[5]{(1) #1 \par (2) #2 \par (3) #3  \par (4) #4  \par (5) #5  }

%Auflistungen

\newcommand{\listitem}{\par \listskip \listdot}

\newcommand{\auflistungzwei}[2]{\par \listskip \listdot  #1 \par \listskip \listdot  #2 \par \listskip}

\newcommand{\auflistungdrei}[3]{\par \listskip \listdot  #1 \par \listskip \listdot  #2\par \listskip \listdot  #3\par\listskip }

\newcommand{\auflistungvier}[4]{\par \listskip \listdot  #1 \par \listskip \listdot  #2\par \listskip \listdot  #3\par \listskip \listdot  #4  \par \listskip }

\newcommand{\auflistungfuenf}[5]{\par \listskip \listdot  #1 \par \listskip \listdot  #2\par \listskip \listdot  #3\par \listskip \listdot  #4  \par \listskip  \listdot  #5 \par \listskip   }

\newcommand{\auflistungsechs}[6]{\par \listskip \listdot  #1 \par \listskip \listdot  #2\par \listskip \listdot  #3\par \listskip \listdot  #4  \par \listskip  \listdot  #5 \par \listskip \listdot  #6 \par \listskip}

\newcommand{\auflistungsieben}[7]{\par \listskip \listdot  #1 \par \listskip \listdot  #2\par \listskip \listdot  #3\par \listskip \listdot  #4  \par \listskip  \listdot  #5 \par \listskip \listdot  #6 \par \listskip \listdot  #7 \par \listskip }

\newcommand{\auflistungacht}[8]{\par \listskip \listdot  #1 \par \listskip \listdot  #2\par \listskip \listdot  #3\par \listskip \listdot  #4  \par \listskip  \listdot  #5 \par \listskip \listdot  #6 \par \listskip \listdot  #7 \par \listskip \listdot  #8  \par \listskip}

\newcommand{\auflistungneun}[9]{\par \listskip \listdot  #1 \par \listskip \listdot  #2\par \listskip \listdot  #3\par \listskip \listdot  #4  \par \listskip  \listdot  #5 \par \listskip \listdot  #6 \par \listskip \listdot  #7 \par \listskip \listdot  #8  \par \listskip  \listdot  #9 \par \listskip  }

\newcommand{\listitemfuenf}[5]{\listitem  #1 \listitem  #2 \listitem  #3 \listitem  #4 \listitem  #5}

\newcommand{\auflistungzweireihe}[2]{ #1 #2 }

%Gestaltung der Auflistungen

\newcommand{\listdot}{$\bullet \,$}

%Bildgestaltung

\newcommand{\inputbild}[1]{\includegraphics[scale]{#1}}

%Tabellengestaltung

%[[Projekt:Semantische Vorlagen/Zeilenkonfiguration/Latex]]

\newcommand{\zeileundzwei}[2]{#1 & #2}

\newcommand{\zeileunddrei}[3]{#1 & #2 & #3}

\newcommand{\zeileundvier}[4]{#1 & #2 & #3 & #4}

\newcommand{\zeileundfuenf}[5]{#1 & #2 & #3 & #4 & #5}

\newcommand{\zeileundsechs}[6]{#1 & #2 & #3 & #4 & #5 & #6}

\newcommand{\zeileundsieben}[7]{#1 & #2 & #3 & #4 & #5 & #6 & #7}

\newcommand{\zeileundacht}[8]{#1 & #2 & #3 & #4 & #5 & #6 & #7 & #8}

\newcommand{\zeileundneun}[9]{#1 & #2 & #3 & #4 & #5 & #6 & #7 & #8 & #9}

\newcommand{\mazeileundzwei}[2]{$#1$ & $#2$}

\newcommand{\mazeileunddrei}[3]{$#1$ & $ #2$ & $#3$}

\newcommand{\mazeileundvier}[4]{$#1$ & $#2$ & $#3$ & $#4$}

\newcommand{\mazeileundfuenf}[5]{$#1$ & $#2$ & $#3$ & $#4$ & $#5$}

\newcommand{\mazeileundsechs}[6]{$#1$ & $#2$ & $#3$ & $#4$ & $#5$ & $#6$}

\newcommand{\mazeileundsieben}[7]{$#1$ & $#2$ & $#3$ & $#4$ & $#5$ & $#6$ & $#7$}

\newcommand{\mazeileundacht}[8]{$#1$ & $#2$ & $#3$ & $#4$ & $#5$ & $#6$ & $#7$ & $#8$}

\newcommand{\mazeileundneun}[9]{$#1$ & $#2$ & $#3$ & $#4$ & $#5$ & $#6$ & $#7$ & $#8$ & $#9$}

%besser als Listen?
%Liste mit und (&) bzw. Bruch.

\newcommand {\listesechsmaund}[6]{$#1$ & $#2$& $#3$ & $#4$ & $#5$ & $#6$}

\newcommand {\listesiebenmaund}[7]{$#1$ & $#2$& $#3$ & $#4$ & $#5$ & $#6$ & $#7$}

\newcommand{\listezweiund}[2]{#1}

\newcommand{\listedreiund}[3]{#1}

\newcommand{\listesiebenund}[7]{#1}

\newcommand{\listesechsbruch}[6]{#1 \\ \hline #2 \\ \hline #3 \\ \hline #4 \\ \hline #5 \\ \hline #6}

\newcommand{\leitzeileunddrei}[3]{#1}

\newcommand{\madoppelzeileunddrei}[6]{$#1$&$#2$&$#3$\\ \hline $#4$&$#5$&$#6$}

%[[Projekt:Semantische Vorlagen/Tabellen/Latex]]

\newcommand{\tabelleviervier}[4]{\par \bigskip \begin{tabular}{|c|c|c|c|} \hline #1 \\ \hline #2 \\ \hline #3 \\ \hline #4 \\ \hline \end{tabular}}

\newcommand{\tabellefuenfvier}[5]{\par \bigskip \begin{tabular}{|c|c|c|c|c|} \hline #1 \\ \hline #2 \\ \hline #3 \\ \hline #4 \\ \hline #5  \\ \hline \end{tabular}}

%Mit math. Einträgen - ohne Leitzeile

\newcommand{\matabellezweivier}[5]{\begin{center} \begin{tabular}{|c|c|} \hline #2 \\ \hline #3 \\ \hline #4 \\ \hline #5  \\ \hline \end{tabular} \end{center} }

\newcommand{\matabellezweifuenf}[6]{\begin{center} \begin{tabular}{|c|c|} \hline     #2 \\ \hline #3 \\ \hline #4 \\ \hline #5 \\ \hline #6  \\ \hline \end{tabular} \end{center} }

\newcommand{\matabellezweisechs}[7]{\begin{center} \begin{tabular}{|c|c|} \hline     #2 \\ \hline #3 \\ \hline #4 \\ \hline #5 \\ \hline #6 \\ \hline #7   \\ \hline \end{tabular} \end{center} }

\newcommand{\matabellezweisieben}[8]{\begin{center} \begin{tabular}{|c|c|} \hline     #2 \\ \hline #3 \\ \hline #4 \\ \hline #5 \\ \hline #6 \\ \hline #7 \\ \hline #8  \\ \hline \end{tabular} \end{center} }

%Mit math. Einträgen - Mit Leitzeile

\newcommand{\matabellesiebenzwoelf}[3]{\begin{center} \begin{tabular}{|c||c|c|c|c|c|c|c|}  \hline  #1   \\ \hline \hline #2 \\ \hline #3 \\ \hline \end{tabular} \end{center} }

\newcommand{\matabellezwoelfzwoelf}[3]{\begin{center} \begin{tabular}{|c||c|c|c|c|c|c|c|c|c|c|c|c|}  \hline  #1   \\ \hline \hline #2 \\ \hline #3 \\ \hline \end{tabular} \end{center} }

\newcommand{\matabelledreineun}[5]{\begin{center} \begin{tabular}{|c|c|c|}\hline  #1 \\ \hline  #2 \\ \hline #3 \\ \hline   #4 \\ \hline  #5 \\ \hline \end{tabular} \end{center} }

\newcommand{\matabelledreizehn}[6]{\begin{center} \begin{tabular}{|c|c|c|}\hline  #1 \\ \hline  #2 \\ \hline #3 \\ \hline   #4 \\ \hline  #5 \\ \hline  #6 \\ \hline \end{tabular} \end{center} }

\newcommand{\matabelledreielf}[6]{\begin{center} \begin{tabular}{|c|c|c|}\hline  #1 \\ \hline  #2 \\ \hline #3 \\ \hline   #4 \\ \hline  #5 \\ \hline  #6 \\ \hline \end{tabular} \end{center} }

%[[Projekt:Semantische Vorlagen/Klausur/Punktetabellen/Latex]]

\newcommand{\punktetabellezehn}
{\vspace{3ex} \begin{center} \renewcommand{\arraystretch}{1.5} \tabcolsep=0.78ex \begin{tabular}{@{}|c|c|c|c|c|c|c|c|c|c|c|c|c|}
\hline Aufgabe & 1 & 2 & 3 & 4 & 5 & 6 & 7 & 8 & 9 & 10 & $\sum$
\\ \hline m\"ogliche Pkt. & \aeins & \azwei & \adrei & \avier & \afuenf & \asechs   & \asieben & \aacht  & \aneun & \azehn & \aelf
\\ \hline erhaltene Pkt.  & \phantom{11} & \phantom{11} & \phantom{11} & \phantom{11} & \phantom{11}& \phantom{11} & \phantom{11} & \phantom{11} & \phantom{11} & & \\ \hline \end{tabular} \end{center} \vspace{2ex} }

\newcommand{\punktetabelleelf}
{\vspace{3ex} \begin{center} \renewcommand{\arraystretch}{1.5} \tabcolsep=0.78ex \begin{tabular}{@{}|c|c|c|c|c|c|c|c|c|c|c|c|c|c|}
\hline Aufgabe & 1 & 2 & 3 & 4 & 5 & 6 & 7 & 8 & 9 & 10 & 11 & $\sum$
\\ \hline m\"ogliche Pkt. & \aeins & \azwei & \adrei & \avier & \afuenf & \asechs   & \asieben & \aacht  & \aneun & \azehn & \aelf & \azwoelf
\\ \hline erhaltene Pkt.  & \phantom{11} & \phantom{11} & \phantom{11} & \phantom{11} & \phantom{11}& \phantom{11} & \phantom{11} & \phantom{11} & \phantom{11} & \phantom{11} & & \\ \hline \end{tabular} \end{center} \vspace{2ex} }

\newcommand{\punktetabellezwoelf}
{\vspace{3ex} \begin{center} \renewcommand{\arraystretch}{1.5} \tabcolsep=0.78ex \begin{tabular}{@{}|c|c|c|c|c|c|c|c|c|c|c|c|c|c|c|}
\hline Aufgabe & 1 & 2 & 3 & 4 & 5 & 6 & 7 & 8 & 9 & 10 & 11 & 12    & $\sum$
\\ \hline m\"ogliche Pkt. & \aeins & \azwei & \adrei & \avier & \afuenf & \asechs   & \asieben & \aacht  & \aneun & \azehn & \aelf & \azwoelf & \adreizehn
\\ \hline erhaltene Pkt.  & \phantom{11} & \phantom{11} & \phantom{11} & \phantom{11} & \phantom{11}& \phantom{11} & \phantom{11} & \phantom{11}
& \phantom{11} & \phantom{11} & \phantom{11} & & \\ \hline \end{tabular} \end{center} \vspace{2ex} }

\newcommand{\punktetabelledreizehn}
{\vspace{3ex} \begin{center} \renewcommand{\arraystretch}{1.5} \tabcolsep=0.78ex \begin{tabular}{@{}|c|c|c|c|c|c|c|c|c|c|c|c|c|c|c|c|}
\hline Aufgabe & 1 & 2 & 3 & 4 & 5 & 6 & 7 & 8 & 9 & 10 & 11 & 12 & 13   & $\sum$
\\ \hline m\"ogliche Pkt. & \aeins & \azwei & \adrei & \avier & \afuenf & \asechs   & \asieben & \aacht  & \aneun & \azehn & \aelf & \azwoelf & \adreizehn & \avierzehn
\\ \hline erhaltene Pkt.  & \phantom{11} & \phantom{11} & \phantom{11} & \phantom{11} & \phantom{11}& \phantom{11} & \phantom{11} & \phantom{11} & \phantom{11}
& \phantom{11} & \phantom{11} & \phantom{11} & & \\ \hline \end{tabular} \end{center} \vspace{2ex} }

\newcommand{\punktetabellevierzehn}
{\vspace{3ex} \begin{center} \renewcommand{\arraystretch}{1.5} \tabcolsep=0.78ex \begin{tabular}{@{}|c|c|c|c|c|c|c|c|c|c|c|c|c|c|c|c|c|}
\hline Aufgabe & 1 & 2 & 3 & 4 & 5 & 6 & 7 & 8 & 9 & 10 & 11 & 12 & 13 & 14  & $\sum$
\\ \hline m\"ogliche Pkt. & \aeins & \azwei & \adrei & \avier & \afuenf & \asechs   & \asieben & \aacht  & \aneun & \azehn & \aelf & \azwoelf & \adreizehn & \avierzehn & \afuenfzehn
\\ \hline erhaltene Pkt. & \phantom{11} & \phantom{11} & \phantom{11} & \phantom{11} & \phantom{11} & \phantom{11}& \phantom{11} & \phantom{11} & \phantom{11} & \phantom{11}
& \phantom{11} & \phantom{11} & \phantom{11} & & \\ \hline \end{tabular} \end{center} \vspace{2ex} }

\newcommand{\punktetabellefuenfzehn}
{\vspace{3ex} \begin{center} \renewcommand{\arraystretch}{1.5} \tabcolsep=0.78ex \begin{tabular}{@{}|c|c|c|c|c|c|c|c|c|c|c|c|c|c|c|c|c|c|}
\hline Aufgabe & 1 & 2 & 3 & 4 & 5 & 6 & 7 & 8 & 9 & 10 & 11 & 12 & 13 & 14 & 15 & $\sum$
\\ \hline m\"ogliche Pkt. & \aeins & \azwei & \adrei & \avier & \afuenf & \asechs   & \asieben & \aacht  & \aneun & \azehn & \aelf & \azwoelf & \adreizehn & \avierzehn & \afuenfzehn
& \asechzehn
\\ \hline erhaltene Pkt. & \phantom{11} & \phantom{11} & \phantom{11}
& \phantom{11} & \phantom{11} & \phantom{11} & \phantom{11}& \phantom{11} & \phantom{11} & \phantom{11} & \phantom{11}
& \phantom{11} & \phantom{11} & \phantom{11} & & \\ \hline \end{tabular} \end{center} \vspace{2ex} }

\newcommand{\punktetabellesechzehn}
{\vspace{3ex} \begin{center} \renewcommand{\arraystretch}{1.5} \tabcolsep=0.78ex \begin{tabular}{@{}|c|c|c|c|c|c|c|c|c|c|c|c|c|c|c|c|c|c|}
\hline Aufgabe & 1 & 2 & 3 & 4 & 5 & 6 & 7 & 8 & 9 & 10 & 11 & 12 & 13 & 14 & 15 & 16 & $\sum$
\\ \hline m\"ogliche Pkt. & \aeins & \azwei & \adrei & \avier & \afuenf & \asechs   & \asieben & \aacht  & \aneun & \azehn & \aelf & \azwoelf & \adreizehn & \avierzehn & \afuenfzehn
& \asechzehn   & \asiebzehn
\\ \hline erhaltene Pkt. & \phantom{11} & \phantom{11} & \phantom{11} & \phantom{11}
& \phantom{11} & \phantom{11} & \phantom{11} & \phantom{11}& \phantom{11} & \phantom{11} & \phantom{11} & \phantom{11}
& \phantom{11} & \phantom{11} & \phantom{11} & & \\ \hline \end{tabular} \end{center} \vspace{2ex} }

\newcommand{\punktetabellesiebzehn}
{\vspace{3ex} \begin{center} \renewcommand{\arraystretch}{1.5} \tabcolsep=0.78ex \begin{tabular}{@{}|c|c|c|c|c|c|c|c|c|c|c|c|c|c|c|c|c|c|c|c|}
\hline Aufgabe & 1 & 2 & 3 & 4 & 5 & 6 & 7 & 8 & 9 & 10 & 11 & 12 & 13 & 14 & 15 & 16 & 17 & $\sum$
\\ \hline m\"ogliche Pkt. & \aeins & \azwei & \adrei & \avier & \afuenf & \asechs   & \asieben & \aacht  & \aneun & \azehn & \aelf & \azwoelf & \adreizehn & \avierzehn & \afuenfzehn
& \asechzehn   & \asiebzehn & \aachtzehn
\\ \hline erhaltene Pkt. & \phantom{11} & \phantom{11} & \phantom{11} & \phantom{11}
& \phantom{11} & \phantom{11} & \phantom{11} & \phantom{11} & \phantom{11}& \phantom{11} & \phantom{11} & \phantom{11} & \phantom{11}
& \phantom{11} & \phantom{11} & \phantom{11} & & \\ \hline \end{tabular} \end{center} \vspace{2ex} }

\newcommand{\punktetabelleachtzehn}
{\vspace{3ex} \begin{center} \renewcommand{\arraystretch}{1.5} \tabcolsep=0.78ex \begin{tabular}{@{}|c|c|c|c|c|c|c|c|c|c|c|c|c|c|c|c|c|c|c|c|c|}
\hline Aufgabe & 1 & 2 & 3 & 4 & 5 & 6 & 7 & 8 & 9 & 10 & 11 & 12 & 13 & 14 & 15 & 16 & 17 & 18 & $\sum$
\\ \hline m\"ogliche Pkt. & \aeins & \azwei & \adrei & \avier & \afuenf & \asechs   & \asieben & \aacht  & \aneun & \azehn & \aelf & \azwoelf & \adreizehn & \avierzehn & \afuenfzehn
& \asechzehn   & \asiebzehn & \aachtzehn & \aneunzehn
\\ \hline erhaltene Pkt. & \phantom{11} & \phantom{11} & \phantom{11} & \phantom{11} & \phantom{11}
& \phantom{11} & \phantom{11} & \phantom{11} & \phantom{11} & \phantom{11}& \phantom{11} & \phantom{11} & \phantom{11} & \phantom{11}
& \phantom{11} & \phantom{11} & \phantom{11} & & \\ \hline \end{tabular} \end{center} \vspace{2ex} }

\newcommand{\punktetabelleneunzehn}
{\vspace{3ex} \begin{center} \renewcommand{\arraystretch}{1.5} \tabcolsep=0.78ex \begin{tabular}{@{}|c|c|c|c|c|c|c|c|c|c|c|c|c|c|c|c|c|c|c|c|c|c|}
\hline Aufgabe & 1 & 2 & 3 & 4 & 5 & 6 & 7 & 8 & 9 & 10 & 11 & 12 & 13 & 14 & 15 & 16 & 17 & 18 & 19 & $\sum$
\\ \hline m\"ogliche Pkt. & \aeins & \azwei & \adrei & \avier & \afuenf & \asechs   & \asieben & \aacht  & \aneun & \azehn & \aelf & \azwoelf & \adreizehn & \avierzehn & \afuenfzehn
& \asechzehn   & \asiebzehn & \aachtzehn & \aneunzehn  & \azwanzig
\\ \hline erhaltene Pkt. & \phantom{11} & \phantom{11} & \phantom{11} & \phantom{11} & \phantom{11} & \phantom{11}
& \phantom{11} & \phantom{11} & \phantom{11} & \phantom{11} & \phantom{11}& \phantom{11} & \phantom{11} & \phantom{11} & \phantom{11}
& \phantom{11} & \phantom{11} & \phantom{11} & & \\ \hline \end{tabular} \end{center} \vspace{2ex} }

\newcommand{\punktetabellezwanzig}
{\vspace{3ex} \begin{center} \renewcommand{\arraystretch}{1.5} \tabcolsep=0.78ex \begin{tabular}{@{}|c|c|c|c|c|c|c|c|c|c|c|c|c|c|c|c|c|c|c|c|c|c|c|}
\hline Aufgabe & 1 & 2 & 3 & 4 & 5 & 6 & 7 & 8 & 9 & 10 & 11 & 12 & 13 & 14 & 15 & 16 & 17 & 18 & 19 & 20 & $\sum$
\\ \hline m\"ogliche Pkt. & \aeins & \azwei & \adrei & \avier & \afuenf & \asechs   & \asieben & \aacht  & \aneun & \azehn & \aelf & \azwoelf & \adreizehn & \avierzehn & \afuenfzehn
& \asechzehn   & \asiebzehn & \aachtzehn & \aneunzehn  & \azwanzig & \aeinundzwanzig
\\ \hline erhaltene Pkt. & \phantom{11} & \phantom{11} & \phantom{11} & \phantom{11} & \phantom{11} & \phantom{11} & \phantom{11}
& \phantom{11} & \phantom{11} & \phantom{11} & \phantom{11} & \phantom{11}& \phantom{11} & \phantom{11} & \phantom{11} & \phantom{11}
& \phantom{11} & \phantom{11} & \phantom{11} & & \\ \hline \end{tabular} \end{center} \vspace{2ex} }

\newcommand{\punktetabelleeinundzwanzig}
{\vspace{3ex} \begin{center} \renewcommand{\arraystretch}{1.5} \tabcolsep=0.78ex \begin{tabular}{@{}|c|c|c|c|c|c|c|c|c|c|c|c|c|c|c|c|c|c|c|c|c|c|c|c|}
\hline Aufgabe & 1 & 2 & 3 & 4 & 5 & 6 & 7 & 8 & 9 & 10 & 11 & 12 & 13 & 14 & 15 & 16 & 17 & 18 & 19 & 20 & 21 & $\sum$
\\ \hline m\"ogliche Pkt. & \aeins & \azwei & \adrei & \avier & \afuenf & \asechs   & \asieben & \aacht  & \aneun & \azehn & \aelf & \azwoelf & \adreizehn & \avierzehn & \afuenfzehn
& \asechzehn   & \asiebzehn & \aachtzehn & \aneunzehn  & \azwanzig & \aeinundzwanzig  & \azweiundzwanzig
\\ \hline erhaltene Pkt. & \phantom{11} & \phantom{11} & \phantom{11} & \phantom{11} & \phantom{11} & \phantom{11} & \phantom{11}
& \phantom{11} & \phantom{11} & \phantom{11} & \phantom{11} & \phantom{11}& \phantom{11} & \phantom{11} & \phantom{11} & \phantom{11}
& \phantom{11} & \phantom{11} & \phantom{11} & \phantom{11} & & \\ \hline \end{tabular} \end{center} \vspace{2ex} }

\newcommand{\punktetabellezweiundzwanzig}
{\vspace{3ex} \begin{center} \renewcommand{\arraystretch}{1.5} \tabcolsep=0.78ex \begin{tabular}{@{}|c|c|c|c|c|c|c|c|c|c|c|c|c|c|c|c|c|c|c|c|c|c|c|c|c|}
\hline Aufgabe & 1 & 2 & 3 & 4 & 5 & 6 & 7 & 8 & 9 & 10 & 11 & 12 & 13 & 14 & 15 & 16 & 17 & 18 & 19 & 20 & 21 & 22 & $\sum$
\\ \hline  Punkte & \aeins & \azwei & \adrei & \avier & \afuenf & \asechs   & \asieben & \aacht  & \aneun & \azehn & \aelf & \azwoelf & \adreizehn & \avierzehn & \afuenfzehn
& \asechzehn   & \asiebzehn & \aachtzehn & \aneunzehn  & \azwanzig & \aeinundzwanzig  & \azweiundzwanzig  & \adreiundzwanzig
\\ \hline Erhalten & \phantom{11} & \phantom{11} & \phantom{11} & \phantom{11} & \phantom{11} & \phantom{11} & \phantom{11}
& \phantom{11} & \phantom{11} & \phantom{11} & \phantom{11} & \phantom{11}& \phantom{11} & \phantom{11} & \phantom{11}  & \phantom{11} & \phantom{11}
& \phantom{11} & \phantom{11} & \phantom{11} & \phantom{11} & & \\ \hline \end{tabular} \end{center} \vspace{2ex} }

%[[Projekt:Semantische Vorlagen/Wahrheitstabellen/Latex]]

\newcommand{\wahrheitstabelleeins}[4]{
\begin{center}
\begin{tabular}{|c|c|}
\multicolumn{2}{l}{\textbf{#1}}\\
\hline
$p$ & $#2$\\ \hline
w  &	#3 \\ \hline
f  &	#4\\ \hline
\end{tabular}
\end{center} }

\newcommand{\wahrheitstabelleeinszwei}[7]{
\begin{center}
\begin{tabular}{|c|c|c|}
\multicolumn{3}{l}{\textbf{#1}}\\
\hline
$p$ & $#2$ & $#5$ \\ \hline
w  &	#3 &  #4 \\ \hline
f  &	#6 &  #7 \\ \hline
\end{tabular}
\end{center} }

\newcommand{\wahrheitstabelleeinsdrei}[4]{ \begin{center} \begin{tabular}{|c|c|c|c|} \multicolumn{4}{l}{\textbf{#1}} \\  \hline #2  \\  \hline #3  \\  \hline  #4  \\  \hline
\end{tabular} \end{center} }

%Wahrheitstabellen mit zwei Variablen p und q

\newcommand{\wahrheitstabellezweieins}[6]{ \begin{center} \begin{tabular}{|c|c|c|} \multicolumn{3}{l}{\textbf{#1}} \\  \hline  #2  \\  \hline  #3  \\  \hline  #4  \\  \hline  #5  \\  \hline  #6  \\  \hline \end{tabular} \end{center} }

\newcommand{\wahrheitstabellezweizwei}[6]{ \begin{center} \begin{tabular}{|c|c|c|c|} \multicolumn{4}{l}{\textbf{#1}} \\  \hline #2  \\  \hline #3  \\  \hline #4  \\  \hline #5  \\  \hline #6  \\  \hline\end{tabular} \end{center} }

\newcommand{\wahrheitstabellezweidrei}[6]{ \begin{center} \begin{tabular}{|c|c|c|c|c|} \multicolumn{5}{l}{\textbf{#1}} \\  \hline #2  \\  \hline #3  \\  \hline #4  \\  \hline #5  \\  \hline #6  \\  \hline \end{tabular} \end{center} }

\newcommand{\wahrheitstabellezweivier}[6]{ \begin{center} \begin{tabular}{|c|c|c|c|c|c|} \multicolumn{6}{l}{\textbf{#1}} \\  \hline #2  \\  \hline #3  \\  \hline #4  \\  \hline  #5  \\  \hline  #6  \\  \hline \end{tabular} \end{center} }

\newcommand{\wahrheitstabellezweifuenf}[6]{ \begin{center} \begin{tabular}{|c|c|c|c|c|c|c|} \multicolumn{7}{l}{\textbf{#1}} \\  \hline #2  \\  \hline #3  \\  \hline #4  \\  \hline #5  \\  \hline #6  \\  \hline
\end{tabular} \end{center} }

%Wahrheitstabellen mit drei Variablen p,q,r

\newcommand{\wahrheitstabelledreieins}[6]{ \begin{center} \begin{tabular}{|c|c|c|c||} \multicolumn{4}{l}{\textbf{#1}} \\  \hline #2  \\  \hline #3  \\  \hline #4  \\  \hline #5  \\  \hline #6 \\  \hline
\end{tabular} \end{center} }

%[[Projekt:Semantische Vorlagen/Wertetabellen/Latex]]

%Wertetabellen

\newcommand{\wertetabellevorskip}{\par \bigskip}

\newcommand{\wertetabellenachskip}{\par \bigskip}

\newcommand{\wertetabelleeinsausteilzeilen}[4]{\wertetabellevorskip \begin{tabular}{|c|c|} \hline #1 & #2  \\ \hline  #3 & #4  \\ \hline \end{tabular} \wertetabellenachskip }

\newcommand{\wertetabellezweiausteilzeilen}[4]{\wertetabellevorskip \begin{tabular}{|c|c|c|} \hline #1 & #2  \\ \hline  #3 & #4  \\ \hline \end{tabular} \wertetabellenachskip }

\newcommand{\wertetabelledreiausteilzeilen}[4]{\wertetabellevorskip \begin{tabular}{|c|c|c|c|} \hline #1 & #2  \\ \hline  #3 & #4  \\ \hline \end{tabular} \wertetabellenachskip }

\newcommand{\wertetabellevierausteilzeilen}[4]{\wertetabellevorskip \begin{tabular}{|c|c|c|c|c|} \hline #1 & #2  \\ \hline  #3 & #4  \\ \hline \end{tabular} \wertetabellenachskip }

\newcommand{\wertetabellefuenfausteilzeilen}[4]{\wertetabellevorskip \begin{tabular}{|c|c|c|c|c|c|} \hline #1 & #2  \\ \hline  #3 & #4  \\ \hline \end{tabular} \wertetabellenachskip }

\newcommand{\wertetabellesechsausteilzeilen}[6]{\wertetabellevorskip \begin{tabular}{|c|c|c|c|c|c|c|} \hline #1 & #2 & #3 \\ \hline  #4 & #5 & #6 \\ \hline \end{tabular} \wertetabellenachskip }

\newcommand{\wertetabellesiebenausteilzeilen}[6]{\wertetabellevorskip \begin{tabular}{|c|c|c|c|c|c|c|c|} \hline #1 & #2 & #3 \\ \hline  #4 & #5 & #6 \\ \hline \end{tabular} \wertetabellenachskip }

\newcommand{\wertetabelleachtausteilzeilen}[6]{\wertetabellevorskip \begin{tabular}{|c|c|c|c|c|c|c|c|c|} \hline #1 & #2 & #3 \\ \hline  #4 & #5 & #6 \\ \hline \end{tabular} \wertetabellenachskip }

\newcommand{\wertetabelleneunausteilzeilen}[6]{\wertetabellevorskip \begin{tabular}{|c|c|c|c|c|c|c|c|c|c|} \hline #1 & #2 & #3 \\ \hline  #4 & #5 & #6 \\ \hline \end{tabular} \wertetabellenachskip }

\newcommand{\wertetabellezehnausteilzeilen}[6]{\wertetabellevorskip \begin{tabular}{|c|c|c|c|c|c|c|c|c|c|c|} \hline #1 & #2 & #3 \\ \hline  #4 & #5 & #6 \\ \hline \end{tabular} \wertetabellenachskip }

\newcommand{\wertetabelleelfausteilzeilen}[8]{\wertetabellevorskip \begin{tabular}{|c|c|c|c|c|c|c|c|c|c|c|c|} \hline #1 & #2 & #3 & #4 \\ \hline #5 & #6 & #7 & #8 \\ \hline \end{tabular} \wertetabellenachskip }

%[[Projekt:Semantische Vorlagen/Verknüpfungstabellen/Latex]]

%Verknüpfungstabellen

\newcommand{\verknuepfungstabelleeins}[2]{\begin{center}
\begin{tabular}{|c||c|}
\hline #1 \\  \hline \hline #2 \\
\hline
\end{tabular}
\end{center} }

\newcommand{\verknuepfungstabellezwei}[3]{\begin{center}
\begin{tabular}{|c||c|c|}
\hline #1 \\  \hline \hline #2 \\ \hline #3 \\
\hline
\end{tabular}
\end{center} }

\newcommand{\verknuepfungstabelledrei}[4]{\begin{center}
\begin{tabular}{|c||c|c|c|}
\hline #1 \\  \hline \hline #2 \\ \hline #3 \\ \hline #4 \\
\hline
\end{tabular}
\end{center} }

\newcommand{\verknuepfungstabellevier}[5]{\begin{center}
\begin{tabular}{|c||c|c|c|c|}
\hline #1 \\  \hline \hline #2 \\ \hline #3 \\ \hline #4 \\  \hline  #5 \\
\hline
\end{tabular}
\end{center} }

\newcommand{\verknuepfungstabellefuenf}[6]{\begin{center}
\begin{tabular}{|c||c|c|c|c|c|}
\hline #1 \\  \hline \hline #2 \\ \hline #3 \\ \hline #4 \\  \hline  #5 \\ \hline  #6 \\
\hline
\end{tabular}
\end{center} }

\newcommand{\verknuepfungstabellesechs}[7]{\begin{center}
\begin{tabular}{|c||c|c|c|c|c|c|}
\hline #1 \\  \hline \hline #2 \\ \hline #3 \\ \hline #4 \\  \hline  #5 \\ \hline  #6 \\  \hline #7 \\
\hline
\end{tabular}
\end{center} }

\newcommand{\verknuepfungstabellesieben}[8]{\begin{center}
\begin{tabular}{|c||c|c|c|c|c|c|c|c|}
\hline #1 \\  \hline \hline #2 \\ \hline #3 \\ \hline #4 \\  \hline  #5 \\ \hline  #6 \\  \hline #7 \\  \hline  #8 \\
\hline
\end{tabular}
\end{center} }

\newcommand{\verknuepfungstabelleacht}[9]{\begin{center}
\begin{tabular}{|c||c|c|c|c|c|c|c|c|}
\hline #1 \\  \hline \hline #2 \\ \hline #3 \\ \hline #4 \\  \hline  #5 \\ \hline  #6 \\  \hline #7 \\  \hline  #8 \\ \hline  #9 \\
\hline
\end{tabular}
\end{center} }

%[[Projekt:Semantische Vorlagen/Initialisierung für Daten/Latex]]

\newcommand{\aeins}{    }

\newcommand{\azwei}{     }

\newcommand{\adrei}{    }

\newcommand{\avier}{    }

\newcommand{\afuenf}{    }

\newcommand{\asechs}{    }

\newcommand{\asieben}{    }

\newcommand{\aacht}{    }

\newcommand{\aneun}{    }

\newcommand{\azehn}{    }

\newcommand{\aelf}{    }

\newcommand{\azwoelf}{    }

\newcommand{\adreizehn}{     }

\newcommand{\avierzehn}{    }

\newcommand{\afuenfzehn}{    }

\newcommand{\asechzehn}{    }

\newcommand{\asiebzehn}{    }

\newcommand{\aachtzehn}{    }

\newcommand{\aneunzehn}{    }

\newcommand{\azwanzig}{    }

\newcommand{\aeinundzwanzig}{    }

\newcommand{\azweiundzwanzig}{    }

\newcommand{\adreiundzwanzig}{    }

\newcommand{\avierundzwanzig}{    }

\newcommand{\afuenfundzwanzig}{    }

\newcommand{\asechsundzwanzig}{    }

\newcommand{\asiebenundzwanzig}{    }

\newcommand{\aachtundzwanzig}{    }

\newcommand{\aneunundzwanzig}{    }

\newcommand{\leitzeilenull}{}

\newcommand{\leitzeileeins}{}

\newcommand{\leitzeilezwei}{}

\newcommand{\leitzeiledrei}{}

\newcommand{\leitzeilevier}{}

\newcommand{\leitzeilefuenf}{}

\newcommand{\leitzeilesechs}{}

\newcommand{\leitzeilesieben}{}

\newcommand{\leitzeileacht}{}

\newcommand{\leitzeileneun}{}

\newcommand{\leitzeilezehn}{}

\newcommand{\leitzeileelf}{}

\newcommand{\leitzeilezwoelf}{}

\newcommand{\leitspaltenull}{}

\newcommand{\leitspalteeins}{}

\newcommand{\leitspaltezwei}{}

\newcommand{\leitspaltedrei}{}

\newcommand{\leitspaltevier}{}

\newcommand{\leitspaltefuenf}{}

\newcommand{\leitspaltesechs}{}

\newcommand{\leitspaltesieben}{}

\newcommand{\leitspalteacht}{}

\newcommand{\leitspalteneun}{}

\newcommand{\leitspaltezehn}{}

\newcommand{\leitspalteelf}{}

\newcommand{\leitspaltezwoelf}{}

\newcommand{\leitspaltedreizehn}{}

\newcommand{\leitspaltevierzehn}{}

\newcommand{\leitspaltefuenfzehn}{}

\newcommand{\leitspaltesechzehn}{}

\newcommand{\leitspaltesiebzehn}{}

\newcommand{\leitspalteachtzehn}{}

\newcommand{\leitspalteneunzehn}{}

\newcommand{\leitspaltezwanzig}{}

\newcommand{\aeinsxeins}{}

\newcommand{\aeinsxzwei}{}

\newcommand{\aeinsxdrei}{}

\newcommand{\aeinsxvier}{}

\newcommand{\aeinsxfuenf}{}

\newcommand{\aeinsxsechs}{}

\newcommand{\aeinsxsieben}{}

\newcommand{\aeinsxacht}{}

\newcommand{\aeinsxneun}{}

\newcommand{\aeinsxzehn}{}

\newcommand{\aeinsxelf}{}

\newcommand{\aeinsxzwoelf}{}

\newcommand{\azweixeins}{}

\newcommand{\azweixzwei}{}

\newcommand{\azweixdrei}{}

\newcommand{\azweixvier}{}

\newcommand{\azweixfuenf}{}

\newcommand{\azweixsechs}{}

\newcommand{\azweixsieben}{}

\newcommand{\azweixacht}{}

\newcommand{\azweixneun}{}

\newcommand{\azweixzehn}{}

\newcommand{\azweixelf}{}

\newcommand{\azweixzwoelf}{}

\newcommand{\adreixeins}{}

\newcommand{\adreixzwei}{}

\newcommand{\adreixdrei}{}

\newcommand{\adreixvier}{}

\newcommand{\adreixfuenf}{}

\newcommand{\adreixsechs}{}

\newcommand{\adreixsieben}{}

\newcommand{\adreixacht}{}

\newcommand{\adreixneun}{}

\newcommand{\adreixzehn}{}

\newcommand{\adreixelf}{}

\newcommand{\adreixzwoelf}{}

\newcommand{\avierxeins}{}

\newcommand{\avierxzwei}{}

\newcommand{\avierxdrei}{}

\newcommand{\avierxvier}{}

\newcommand{\avierxfuenf}{}

\newcommand{\avierxsechs}{}

\newcommand{\avierxsieben}{}

\newcommand{\avierxacht}{}

\newcommand{\avierxneun}{}

\newcommand{\avierxzehn}{}

\newcommand{\avierxelf}{}

\newcommand{\avierxzwoelf}{}

\newcommand{\afuenfxeins}{}

\newcommand{\afuenfxzwei}{}

\newcommand{\afuenfxdrei}{}

\newcommand{\afuenfxvier}{}

\newcommand{\afuenfxfuenf}{}

\newcommand{\afuenfxsechs}{}

\newcommand{\afuenfxsieben}{}

\newcommand{\afuenfxacht}{}

\newcommand{\afuenfxneun}{}

\newcommand{\afuenfxzehn}{}

\newcommand{\afuenfxelf}{}

\newcommand{\afuenfxzwoelf}{}

\newcommand{\asechsxeins}{}

\newcommand{\asechsxzwei}{}

\newcommand{\asechsxdrei}{}

\newcommand{\asechsxvier}{}

\newcommand{\asechsxfuenf}{}

\newcommand{\asechsxsechs}{}

\newcommand{\asechsxsieben}{}

\newcommand{\asechsxacht}{}

\newcommand{\asechsxneun}{}

\newcommand{\asechsxzehn}{}

\newcommand{\asechsxelf}{}

\newcommand{\asechsxzwoelf}{}

\newcommand{\asiebenxeins}{}

\newcommand{\asiebenxzwei}{}

\newcommand{\asiebenxdrei}{}

\newcommand{\asiebenxvier}{}

\newcommand{\asiebenxfuenf}{}

\newcommand{\asiebenxsechs}{}

\newcommand{\asiebenxsieben}{}

\newcommand{\asiebenxacht}{}

\newcommand{\asiebenxneun}{}

\newcommand{\asiebenxzehn}{}

\newcommand{\asiebenxelf}{}

\newcommand{\asiebenxzwoelf}{}

\newcommand{\aachtxeins}{}

\newcommand{\aachtxzwei}{}

\newcommand{\aachtxdrei}{}

\newcommand{\aachtxvier}{}

\newcommand{\aachtxfuenf}{}

\newcommand{\aachtxsechs}{}

\newcommand{\aachtxsieben}{}

\newcommand{\aachtxacht}{}

\newcommand{\aachtxneun}{}

\newcommand{\aachtxzehn}{}

\newcommand{\aachtxelf}{}

\newcommand{\aachtxzwoelf}{}

\newcommand{\aneunxeins}{}

\newcommand{\aneunxzwei}{}

\newcommand{\aneunxdrei}{}

\newcommand{\aneunxvier}{}

\newcommand{\aneunxfuenf}{}

\newcommand{\aneunxsechs}{}

\newcommand{\aneunxsieben}{}

\newcommand{\aneunxacht}{}

\newcommand{\aneunxneun}{}

\newcommand{\aneunxzehn}{}

\newcommand{\aneunxelf}{}

\newcommand{\aneunxzwoelf}{}

\newcommand{\azehnxeins}{}

\newcommand{\azehnxzwei}{}

\newcommand{\azehnxdrei}{}

\newcommand{\azehnxvier}{}

\newcommand{\azehnxfuenf}{}

\newcommand{\azehnxsechs}{}

\newcommand{\azehnxsieben}{}

\newcommand{\azehnxacht}{}

\newcommand{\azehnxneun}{}

\newcommand{\azehnxzehn}{}

\newcommand{\azehnxelf}{}

\newcommand{\azehnxzwoelf}{}

\newcommand{\aelfxeins}{}

\newcommand{\aelfxzwei}{}

\newcommand{\aelfxdrei}{}

\newcommand{\aelfxvier}{}

\newcommand{\aelfxfuenf}{}

\newcommand{\aelfxsechs}{}

\newcommand{\aelfxsieben}{}

\newcommand{\aelfxacht}{}

\newcommand{\aelfxneun}{}

\newcommand{\aelfxzehn}{}

\newcommand{\aelfxelf}{}

\newcommand{\aelfxzwoelf}{}

\newcommand{\azwoelfxeins}{}

\newcommand{\azwoelfxzwei}{}

\newcommand{\azwoelfxdrei}{}

\newcommand{\azwoelfxvier}{}

\newcommand{\azwoelfxfuenf}{}

\newcommand{\azwoelfxsechs}{}

\newcommand{\azwoelfxsieben}{}

\newcommand{\azwoelfxacht}{}

\newcommand{\azwoelfxneun}{}

\newcommand{\azwoelfxzehn}{}

\newcommand{\azwoelfxelf}{}

\newcommand{\azwoelfxzwoelf}{}

\newcommand{\adreizehnxeins}{}

\newcommand{\adreizehnxzwei}{}

\newcommand{\adreizehnxdrei}{}

\newcommand{\adreizehnxvier}{}

\newcommand{\adreizehnxfuenf}{}

\newcommand{\adreizehnxsechs}{}

\newcommand{\adreizehnxsieben}{}

\newcommand{\adreizehnxacht}{}

\newcommand{\adreizehnxneun}{}

\newcommand{\adreizehnxzehn}{}

\newcommand{\adreizehnxelf}{}

\newcommand{\adreizehnxzwoelf}{}

\newcommand{\avierzehnxeins}{}

\newcommand{\avierzehnxzwei}{}

\newcommand{\avierzehnxdrei}{}

\newcommand{\avierzehnxvier}{}

\newcommand{\avierzehnxfuenf}{}

\newcommand{\avierzehnxsechs}{}

\newcommand{\avierzehnxsieben}{}

\newcommand{\avierzehnxacht}{}

\newcommand{\avierzehnxneun}{}

\newcommand{\avierzehnxzehn}{}

\newcommand{\avierzehnxelf}{}

\newcommand{\avierzehnxzwoelf}{}

\newcommand{\afuenfzehnxeins}{}

\newcommand{\afuenfzehnxzwei}{}

\newcommand{\afuenfzehnxdrei}{}

\newcommand{\afuenfzehnxvier}{}

\newcommand{\afuenfzehnxfuenf}{}

\newcommand{\afuenfzehnxsechs}{}

\newcommand{\afuenfzehnxsieben}{}

\newcommand{\afuenfzehnxacht}{}

\newcommand{\afuenfzehnxneun}{}

\newcommand{\afuenfzehnxzehn}{}

\newcommand{\afuenfzehnxelf}{}

\newcommand{\afuenfzehnxzwoelf}{}

\newcommand{\asechzehnxeins}{}

\newcommand{\asechzehnxzwei}{}

\newcommand{\asechzehnxdrei}{}

\newcommand{\asechzehnxvier}{}

\newcommand{\asechzehnxfuenf}{}

\newcommand{\asechzehnxsechs}{}

\newcommand{\asechzehnxsieben}{}

\newcommand{\asechzehnxacht}{}

\newcommand{\asechzehnxneun}{}

\newcommand{\asechzehnxzehn}{}

\newcommand{\asechzehnxelf}{}

\newcommand{\asechzehnxzwoelf}{}

\newcommand{\asiebzehnxeins}{}

\newcommand{\asiebzehnxzwei}{}

\newcommand{\asiebzehnxdrei}{}

\newcommand{\asiebzehnxvier}{}

\newcommand{\asiebzehnxfuenf}{}

\newcommand{\asiebzehnxsechs}{}

\newcommand{\asiebzehnxsieben}{}

\newcommand{\asiebzehnxacht}{}

\newcommand{\asiebzehnxneun}{}

\newcommand{\asiebzehnxzehn}{}

\newcommand{\asiebzehnxelf}{}

\newcommand{\asiebzehnxzwoelf}{}

\newcommand{\aachtzehnxeins}{}

\newcommand{\aachtzehnxzwei}{}

\newcommand{\aachtzehnxdrei}{}

\newcommand{\aachtzehnxvier}{}

\newcommand{\aachtzehnxfuenf}{}

\newcommand{\aachtzehnxsechs}{}

\newcommand{\aachtzehnxsieben}{}

\newcommand{\aachtzehnxacht}{}

\newcommand{\aachtzehnxneun}{}

\newcommand{\aachtzehnxzehn}{}

\newcommand{\aachtzehnxelf}{}

\newcommand{\aachtzehnxzwoelf}{}

%[[Projekt:Semantische Vorlagen/Tabellen mit benannten Parametern/Latex]]

%Tabellen mit zwei Zeilen

\newcommand{\tabelleleitzweixzwei}{

\begin{center}

\begin{tabular}{|c||c|c|}

\hline  \leitzeilenull  &  \leitzeileeins  &  \leitzeilezwei   \\

\hline \hline  \leitspalteeins  &  $ \aeinsxeins $  &  $ \aeinsxzwei $ \\

\hline  \leitspaltezwei  &  $ \azweixeins $  &  $ \azweixzwei $   \\

\hline

\end{tabular}

\end{center} }

%Tabellen mit drei Zeilen

\newcommand{\tabelleleitdreixdrei}{

\begin{center}

\begin{tabular}{|c||c|c|c||}

\hline  \leitzeilenull  &  \leitzeileeins  &  \leitzeilezwei  &  \leitzeiledrei        \\

\hline \hline  \leitspalteeins  &  $ \aeinsxeins $  &  $ \aeinsxzwei $  &  $ \aeinsxdrei $      \\

\hline  \leitspaltezwei  &  $ \azweixeins $  &  $ \azweixzwei $  &  $ \azweixdrei $       \\

\hline  \leitspaltedrei  &  $ \adreixeins $  &  $ \adreixzwei $  &  $ \adreixdrei $       \\

\hline

\end{tabular}

\end{center} }

\newcommand{\tabelleleitdreixvier}{

\begin{center}

\begin{tabular}{|c||c|c|c|c||}

\hline  \leitzeilenull  &  \leitzeileeins  &  \leitzeilezwei  &  \leitzeiledrei  &  \leitzeilevier      \\

\hline \hline  \leitspalteeins  &  $ \aeinsxeins $  &  $ \aeinsxzwei $  &  $ \aeinsxdrei $  &  $ \aeinsxvier $   \\

\hline  \leitspaltezwei  &  $ \azweixeins $  &  $ \azweixzwei $  &  $ \azweixdrei $  &  $ \azweixvier $     \\

\hline  \leitspaltedrei  &  $ \adreixeins $  &  $ \adreixzwei $  &  $ \adreixdrei $  &  $ \adreixvier $     \\

\hline

\end{tabular}

\end{center} }

\newcommand{\tabelleleitdreixsechs}{

\begin{center}

\begin{tabular}{|c||c|c|c|c|c|c|}

\hline  \leitzeilenull  &  \leitzeileeins  &  \leitzeilezwei  &  \leitzeiledrei  &  \leitzeilevier  &  \leitzeilefuenf   &  \leitzeilesechs    \\

\hline \hline  \leitspalteeins  &  $ \aeinsxeins $  &  $ \aeinsxzwei $  &  $ \aeinsxdrei $  &  $ \aeinsxvier $  &  $ \aeinsxfuenf $  &  $ \aeinsxsechs $  \\

\hline  \leitspaltezwei  &  $ \azweixeins $  &  $ \azweixzwei $  &  $ \azweixdrei $  &  $ \azweixvier $  &  $ \azweixfuenf $   &  $ \azweixsechs $  \\

\hline  \leitspaltedrei  &  $ \adreixeins $  &  $ \adreixzwei $  &  $ \adreixdrei $  &  $ \adreixvier $  &  $ \adreixfuenf $   &  $ \adreixsechs $   \\

\hline

\end{tabular}

\end{center} }

\newcommand{\tabelleleitdreixelf}{

\begin{center}

\begin{tabular}{|c||c|c|c|c|c|c|c|c|c|c|c|}

\hline  \leitzeilenull  &  \leitzeileeins  &  \leitzeilezwei  &  \leitzeiledrei  &  \leitzeilevier  &  \leitzeilefuenf   &  \leitzeilesechs  &  \leitzeilesieben  &  \leitzeileacht  &  \leitzeileneun  &  \leitzeilezehn &  \leitzeileelf    \\

\hline \hline  \leitspalteeins  &  $ \aeinsxeins $  &  $ \aeinsxzwei $  &  $ \aeinsxdrei $  &  $ \aeinsxvier $  &  $ \aeinsxfuenf $  &  $ \aeinsxsechs $  &  $ \aeinsxsieben $  &  $ \aeinsxacht $  &  $ \aeinsxneun $ &  $ \aeinsxzehn $  &  $ \aeinsxelf $  \\

\hline  \leitspaltezwei  &  $ \azweixeins $  &  $ \azweixzwei $  &  $ \azweixdrei $  &  $ \azweixvier $  &  $ \azweixfuenf $   &  $ \azweixsechs $  &  $ \azweixsieben $  &  $ \azweixacht $  &  $ \azweixneun $  &  $ \azweixzehn $ &  $ \azweixelf $     \\

\hline  \leitspaltedrei  &  $ \adreixeins $  &  $ \adreixzwei $  &  $ \adreixdrei $  &  $ \adreixvier $  &  $ \adreixfuenf $   &  $ \adreixsechs $  &  $ \adreixsieben $  &  $ \adreixacht $  &  $ \adreixneun $  &  $ \adreixzehn $     &  $ \adreixelf $   \\

\hline

\end{tabular}

\end{center} }

%Tabellen mit vier Zeilen

\newcommand{\tabelleleitvierxzwei}{

\begin{center}

\begin{tabular}{|c||c|c|}

\hline  \leitzeilenull  &  \leitzeileeins  &  \leitzeilezwei   \\

\hline \hline  \leitspalteeins  &  $ \aeinsxeins $  &  $ \aeinsxzwei $    \\

\hline  \leitspaltezwei  &  $ \azweixeins $  &  $ \azweixzwei $    \\

\hline  \leitspaltedrei  &  $ \adreixeins $  &  $ \adreixzwei $    \\

\hline  \leitspaltevier  &  $ \avierxeins $  &  $ \avierxzwei $    \\

\hline

\end{tabular}

\end{center} }

\newcommand{\tabelleleitvierxdrei}{

\begin{center}

\begin{tabular}{|c||c|c|c|}

\hline  \leitzeilenull  &  \leitzeileeins  &  \leitzeilezwei  &  \leitzeiledrei     \\

\hline \hline  \leitspalteeins  &  $ \aeinsxeins $  &  $ \aeinsxzwei $  &  $ \aeinsxdrei $     \\

\hline  \leitspaltezwei  &  $ \azweixeins $  &  $ \azweixzwei $  &  $ \azweixdrei $      \\

\hline  \leitspaltedrei  &  $ \adreixeins $  &  $ \adreixzwei $  &  $ \adreixdrei $    \\

\hline  \leitspaltevier  &  $ \avierxeins $  &  $ \avierxzwei $  &  $ \avierxdrei $    \\

\hline

\end{tabular}

\end{center} }

\newcommand{\tabelleleitvierxvier}{

\begin{center}

\begin{tabular}{|c||c|c|c|c|}

\hline  \leitzeilenull  &  \leitzeileeins  &  \leitzeilezwei  &  \leitzeiledrei  &  \leitzeilevier    \\

\hline \hline  \leitspalteeins  &  $ \aeinsxeins $  &  $ \aeinsxzwei $  &  $ \aeinsxdrei $  &  $ \aeinsxvier $   \\

\hline  \leitspaltezwei  &  $ \azweixeins $  &  $ \azweixzwei $  &  $ \azweixdrei $  &  $ \azweixvier $     \\

\hline  \leitspaltedrei  &  $ \adreixeins $  &  $ \adreixzwei $  &  $ \adreixdrei $  &  $ \adreixvier $   \\

\hline  \leitspaltevier  &  $ \avierxeins $  &  $ \avierxzwei $  &  $ \avierxdrei $  &  $ \avierxvier $   \\

\hline

\end{tabular}

\end{center} }

\newcommand{\tabelleleittextvierxzwei}{

	\begin{center}

		\begin{tabular}{|c||c|c|}

			\hline  \leitzeilenull  &  \leitzeileeins  &  \leitzeilezwei   \\

			\hline \hline  \leitspalteeins  &   \aeinsxeins   &   \aeinsxzwei    \\

			\hline  \leitspaltezwei  &  \azweixeins   &   \azweixzwei     \\

			\hline  \leitspaltedrei  &  \adreixeins   &   \adreixzwei     \\

			\hline  \leitspaltevier  &  \avierxeins   &   \avierxzwei     \\

			\hline

		\end{tabular}

	\end{center} }

\newcommand{\tabelleleitvierxacht}{

\begin{center}

\begin{tabular}{|c||c|c|c|c|c|c|c|c|}

\hline  \leitzeilenull  &  \leitzeileeins  &  \leitzeilezwei  &  \leitzeiledrei  &  \leitzeilevier  &  \leitzeilefuenf  &  \leitzeilesechs  &  \leitzeilesieben  &  \leitzeileacht  \\

\hline \hline  \leitspalteeins  &  $ \aeinsxeins $  &  $ \aeinsxzwei $  &  $ \aeinsxdrei $  &  $ \aeinsxvier $  &  $ \aeinsxfuenf $  &  $ \aeinsxsechs $  &  $ \aeinsxsieben $  &  $ \aeinsxacht $  \\

\hline  \leitspaltezwei  &  $ \azweixeins $  &  $ \azweixzwei $  &  $ \azweixdrei $  &  $ \azweixvier $   &  $ \azweixfuenf $  &  $ \azweixsechs $  &  $ \azweixsieben $  &  $ \azweixacht $   \\

\hline  \leitspaltedrei  &  $ \adreixeins $  &  $ \adreixzwei $  &  $ \adreixdrei $  &  $ \adreixvier $ &  $ \adreixfuenf $  &  $ \adreixsechs $  &  $ \adreixsieben $  &  $ \adreixacht $   \\

\hline  \leitspaltevier  &  $ \avierxeins $  &  $ \avierxzwei $  &  $ \avierxdrei $  &  $ \avierxvier $  &  $ \avierxfuenf $  &  $ \avierxsechs $  &  $ \avierxsieben $  &  $ \avierxacht $  \\

\hline

\end{tabular}

\end{center} }

%Tabellen mit fünf Zeilen

\newcommand{\tabelleleitfuenfxfuenf}{

\begin{center}

\begin{tabular}{|c||c|c|c|c|c|}

\hline  \leitzeilenull  &  \leitzeileeins  &  \leitzeilezwei  &  \leitzeiledrei  &  \leitzeilevier  &  \leitzeilefuenf   \\

\hline \hline  \leitspalteeins  &  $ \aeinsxeins $  &  $ \aeinsxzwei $  &  $ \aeinsxdrei $  &  $ \aeinsxvier $  &  $ \aeinsxfuenf $   \\

\hline  \leitspaltezwei  &  $ \azweixeins $  &  $ \azweixzwei $  &  $ \azweixdrei $  &  $ \azweixvier $  &  $ \azweixfuenf $     \\

\hline  \leitspaltedrei  &  $ \adreixeins $  &  $ \adreixzwei $  &  $ \adreixdrei $  &  $ \adreixvier $  &  $ \adreixfuenf $      \\

\hline  \leitspaltevier  &  $ \avierxeins $  &  $ \avierxzwei $  &  $ \avierxdrei $  &  $ \avierxvier $  &  $ \avierxfuenf $     \\

\hline  \leitspaltefuenf  &  $ \afuenfxeins $  &  $ \afuenfxzwei $  &  $ \afuenfxdrei $  &  $ \afuenfxvier $  &  $ \afuenfxfuenf $    \\

\hline

\end{tabular}

\end{center} }

%Tabellen mit sechs Zeilen

\newcommand{\tabelleleitsechsxdrei}{

\begin{center}

\begin{tabular}{|c||c|c|c|}

\hline  \leitzeilenull  &  \leitzeileeins  &  \leitzeilezwei  &  \leitzeiledrei    \\

\hline \hline  \leitspalteeins  &  $ \aeinsxeins $  &  $ \aeinsxzwei $  &  $ \aeinsxdrei $    \\

\hline  \leitspaltezwei  &  $ \azweixeins $  &  $ \azweixzwei $  &  $ \azweixdrei $      \\

\hline  \leitspaltedrei  &  $ \adreixeins $  &  $ \adreixzwei $  &  $ \adreixdrei $     \\

\hline  \leitspaltevier  &  $ \avierxeins $  &  $ \avierxzwei $  &  $ \avierxdrei $      \\

\hline  \leitspaltefuenf  &  $ \afuenfxeins $  &  $ \afuenfxzwei $  &  $ \afuenfxdrei $     \\

\hline  \leitspaltesechs  &  $ \asechsxeins $  &  $ \asechsxzwei $  &  $ \asechsxdrei $      \\

\hline

\end{tabular}

\end{center} }

\newcommand{\tabelleleitsechsxfuenf}{

\begin{center}

\begin{tabular}{|c||c|c|c|c|c|}

\hline  \leitzeilenull  &  \leitzeileeins  &  \leitzeilezwei  &  \leitzeiledrei  &  \leitzeilevier  &  \leitzeilefuenf   \\

\hline \hline  \leitspalteeins  &  $ \aeinsxeins $  &  $ \aeinsxzwei $  &  $ \aeinsxdrei $  &  $ \aeinsxvier $  &  $ \aeinsxfuenf $   \\

\hline  \leitspaltezwei  &  $ \azweixeins $  &  $ \azweixzwei $  &  $ \azweixdrei $  &  $ \azweixvier $  &  $ \azweixfuenf $     \\

\hline  \leitspaltedrei  &  $ \adreixeins $  &  $ \adreixzwei $  &  $ \adreixdrei $  &  $ \adreixvier $  &  $ \adreixfuenf $      \\

\hline  \leitspaltevier  &  $ \avierxeins $  &  $ \avierxzwei $  &  $ \avierxdrei $  &  $ \avierxvier $  &  $ \avierxfuenf $     \\

\hline  \leitspaltefuenf  &  $ \afuenfxeins $  &  $ \afuenfxzwei $  &  $ \afuenfxdrei $  &  $ \afuenfxvier $  &  $ \afuenfxfuenf $    \\

\hline  \leitspaltesechs  &  $ \asechsxeins $  &  $ \asechsxzwei $  &  $ \asechsxdrei $  &  $ \asechsxvier $  &  $ \asechsxfuenf $    \\

\hline

\end{tabular}

\end{center} }

\newcommand{\tabelleleitsechsxsechs}{

\begin{center}

\begin{tabular}{|c||c|c|c|c|c|c|}

\hline  \leitzeilenull  &  \leitzeileeins  &  \leitzeilezwei  &  \leitzeiledrei  &  \leitzeilevier  &  \leitzeilefuenf  &  \leitzeilesechs   \\

\hline \hline  \leitspalteeins  &  $ \aeinsxeins $  &  $ \aeinsxzwei $  &  $ \aeinsxdrei $  &  $ \aeinsxvier $  &  $ \aeinsxfuenf $  &  $ \aeinsxsechs $ \\

\hline  \leitspaltezwei  &  $ \azweixeins $  &  $ \azweixzwei $  &  $ \azweixdrei $  &  $ \azweixvier $  &  $ \azweixfuenf $  &  $ \azweixsechs $   \\

\hline  \leitspaltedrei  &  $ \adreixeins $  &  $ \adreixzwei $  &  $ \adreixdrei $  &  $ \adreixvier $  &  $ \adreixfuenf $   &  $ \adreixsechs $   \\

\hline  \leitspaltevier  &  $ \avierxeins $  &  $ \avierxzwei $  &  $ \avierxdrei $  &  $ \avierxvier $  &  $ \avierxfuenf $   &  $ \avierxsechs $  \\

\hline  \leitspaltefuenf  &  $ \afuenfxeins $  &  $ \afuenfxzwei $  &  $ \afuenfxdrei $  &  $ \afuenfxvier $  &  $ \afuenfxfuenf $ &  $ \afuenfxsechs $   \\

\hline  \leitspaltesechs  &  $ \asechsxeins $  &  $ \asechsxzwei $  &  $ \asechsxdrei $  &  $ \asechsxvier $  &  $ \asechsxfuenf $  &  $ \asechsxsechs $  \\

\hline

\end{tabular}

\end{center} }

%Tabellen mit sieben Zeilen

\newcommand{\tabelleleitsiebenxsieben}{

\begin{center}

\begin{tabular}{|c||c|c|c|c|c|c|c|}

\hline  \leitzeilenull  &  \leitzeileeins  &  \leitzeilezwei  &  \leitzeiledrei  &  \leitzeilevier  &  \leitzeilefuenf  &  \leitzeilesechs &  \leitzeilesieben  \\

\hline \hline  \leitspalteeins  &  $ \aeinsxeins $  &  $ \aeinsxzwei $  &  $ \aeinsxdrei $  &  $ \aeinsxvier $  &  $ \aeinsxfuenf $  &  $ \aeinsxsechs $ &  $ \aeinsxsieben $ \\

\hline  \leitspaltezwei  &  $ \azweixeins $  &  $ \azweixzwei $  &  $ \azweixdrei $  &  $ \azweixvier $  &  $ \azweixfuenf $  &  $ \azweixsechs $  &  $ \azweixsieben $ \\

\hline  \leitspaltedrei  &  $ \adreixeins $  &  $ \adreixzwei $  &  $ \adreixdrei $  &  $ \adreixvier $  &  $ \adreixfuenf $   &  $ \adreixsechs $ &  $ \adreixsieben $  \\

\hline  \leitspaltevier  &  $ \avierxeins $  &  $ \avierxzwei $  &  $ \avierxdrei $  &  $ \avierxvier $  &  $ \avierxfuenf $   &  $ \avierxsechs $ &  $ \avierxsieben $ \\

\hline  \leitspaltefuenf  &  $ \afuenfxeins $  &  $ \afuenfxzwei $  &  $ \afuenfxdrei $  &  $ \afuenfxvier $  &  $ \afuenfxfuenf $ &  $ \afuenfxsechs $ &  $ \afuenfxsieben $  \\

\hline  \leitspaltesechs  &  $ \asechsxeins $  &  $ \asechsxzwei $  &  $ \asechsxdrei $  &  $ \asechsxvier $  &  $ \asechsxfuenf $  &  $ \asechsxsechs $ &  $ \asechsxsieben $ \\

\hline  \leitspaltesieben  &  $ \asiebenxeins $  &  $ \asiebenxzwei $  &  $ \asiebenxdrei $  &  $ \asiebenxvier $  &  $ \asiebenxfuenf $  &  $ \asiebenxsechs $ &  $ \asiebenxsieben $ \\

\hline

\end{tabular}

\end{center} }

%Tabellen mit acht Zeilen

\newcommand{\tabelleleitachtxdrei}{

\begin{center}

\begin{tabular}{|c||c|c|c|}

\hline  \leitzeilenull  &  \leitzeileeins  &  \leitzeilezwei  &  \leitzeiledrei     \\

\hline \hline  \leitspalteeins  &  $ \aeinsxeins $  &  $ \aeinsxzwei $  &  $ \aeinsxdrei $    \\

\hline  \leitspaltezwei  &  $ \azweixeins $  &  $ \azweixzwei $  &  $ \azweixdrei $     \\

\hline  \leitspaltedrei  &  $ \adreixeins $  &  $ \adreixzwei $  &  $ \adreixdrei $    \\

\hline  \leitspaltevier  &  $ \avierxeins $  &  $ \avierxzwei $  &  $ \avierxdrei $      \\

\hline  \leitspaltefuenf  &  $ \afuenfxeins $  &  $ \afuenfxzwei $  &  $ \afuenfxdrei $     \\

\hline  \leitspaltesechs  &  $ \asechsxeins $  &  $ \asechsxzwei $  &  $ \asechsxdrei $      \\

\hline  \leitspaltesieben  &  $ \asiebenxeins $  &  $ \asiebenxzwei $  &  $ \asiebenxdrei $     \\

\hline  \leitspalteacht  &  $ \aachtxeins $  &  $ \aachtxzwei $  &  $ \aachtxdrei $      \\

\hline

\end{tabular}

\end{center} }

%Tabellen mit zehn Zeilen

\newcommand{\tabelleleitzehnxdrei}{

\begin{center}

\begin{tabular}{|c||c|c|c|}

\hline  \leitzeilenull  &  \leitzeileeins  &  \leitzeilezwei  &  \leitzeiledrei     \\

\hline \hline  \leitspalteeins  &  $ \aeinsxeins $  &  $ \aeinsxzwei $  &  $ \aeinsxdrei $    \\

\hline  \leitspaltezwei  &  $ \azweixeins $  &  $ \azweixzwei $  &  $ \azweixdrei $     \\

\hline  \leitspaltedrei  &  $ \adreixeins $  &  $ \adreixzwei $  &  $ \adreixdrei $    \\

\hline  \leitspaltevier  &  $ \avierxeins $  &  $ \avierxzwei $  &  $ \avierxdrei $      \\

\hline  \leitspaltefuenf  &  $ \afuenfxeins $  &  $ \afuenfxzwei $  &  $ \afuenfxdrei $     \\

\hline  \leitspaltesechs  &  $ \asechsxeins $  &  $ \asechsxzwei $  &  $ \asechsxdrei $      \\

\hline  \leitspaltesieben  &  $ \asiebenxeins $  &  $ \asiebenxzwei $  &  $ \asiebenxdrei $     \\

\hline  \leitspalteacht  &  $ \aachtxeins $  &  $ \aachtxzwei $  &  $ \aachtxdrei $      \\

\hline  \leitspalteneun  &  $ \aneunxeins $  &  $ \aneunxzwei $  &  $ \aneunxdrei $      \\

\hline  \leitspaltezehn  &  $ \azehnxeins $  &  $ \azehnxzwei $  &  $ \azehnxdrei $      \\

\hline

\end{tabular}

\end{center} }

%Tabellen mit elf Zeilen

\newcommand{\tabelleleitelfxvier}{

\begin{center}

\begin{tabular}{|c||c|c|c|c|}

\hline  \leitzeilenull  &  \leitzeileeins  &  \leitzeilezwei  &  \leitzeiledrei  &  \leitzeilevier  \\

\hline \hline  \leitspalteeins  &  $ \aeinsxeins $  &  $ \aeinsxzwei $  &  $ \aeinsxdrei $  &  $ \aeinsxvier $  \\

\hline  \leitspaltezwei  &  $ \azweixeins $  &  $ \azweixzwei $  &  $ \azweixdrei $  &  $ \azweixvier $   \\

\hline  \leitspaltedrei  &  $ \adreixeins $  &  $ \adreixzwei $  &  $ \adreixdrei $  &  $ \adreixvier $   \\

\hline  \leitspaltevier  &  $ \avierxeins $  &  $ \avierxzwei $  &  $ \avierxdrei $  &  $ \avierxvier $   \\

\hline  \leitspaltefuenf  &  $ \afuenfxeins $  &  $ \afuenfxzwei $  &  $ \afuenfxdrei $  &  $ \afuenfxvier $   \\

\hline  \leitspaltesechs  &  $ \asechsxeins $  &  $ \asechsxzwei $  &  $ \asechsxdrei $  &  $ \asechsxvier $   \\

\hline  \leitspaltesieben  &  $ \asiebenxeins $  &  $ \asiebenxzwei $  &  $ \asiebenxdrei $  &  $ \asiebenxvier $    \\

\hline  \leitspalteacht  &  $ \aachtxeins $  &  $ \aachtxzwei $  &  $ \aachtxdrei $  &  $ \aachtxvier $   \\

\hline  \leitspalteneun  &  $ \aneunxeins $  &  $ \aneunxzwei $  &  $ \aneunxdrei $  &  $ \aneunxvier $   \\

\hline  \leitspaltezehn  &  $ \azehnxeins $  &  $ \azehnxzwei $  &  $ \azehnxdrei $  &  $ \azehnxvier $   \\

\hline  \leitspalteelf  &  $ \aelfxeins $  &  $ \aelfxzwei $  &  $ \aelfxdrei $  &  $ \aelfxvier $   \\

\hline

\end{tabular}

\end{center} }

%Tabellen mit zwölf Zeilen

\newcommand{\tabelleleitzwoelfxsieben}{

\begin{center}

\begin{tabular}{|c||c|c|c|c|c|c|}

\hline  \leitzeilenull  &  \leitzeileeins  &  \leitzeilezwei  &  \leitzeiledrei  &  \leitzeilevier &  \leitzeilefuenf  &  \leitzeilesechs &  \leitzeilesieben \\

\hline \hline  \leitspalteeins  &  $ \aeinsxeins $  &  $ \aeinsxzwei $  &  $ \aeinsxdrei $  &  $ \aeinsxvier $  &  $ \aeinsxfuenf $  &  $ \aeinsxsechs $  &  $ \aeinsxsieben $  \\

\hline  \leitspaltezwei  &  $ \azweixeins $  &  $ \azweixzwei $  &  $ \azweixdrei $  &  $ \azweixvier $  &  $ \azweixfuenf $  &  $ \azweixsechs $  &  $ \azweixsieben $   \\

\hline  \leitspaltedrei  &  $ \adreixeins $  &  $ \adreixzwei $  &  $ \adreixdrei $  &  $ \adreixvier $  &  $ \adreixfuenf $  &  $ \adreixsechs $  &  $ \adreixsieben $   \\

\hline  \leitspaltevier  &  $ \avierxeins $  &  $ \avierxzwei $  &  $ \avierxdrei $  &  $ \avierxvier $   &  $ \avierxfuenf $  &  $ \avierxsechs $  &  $ \avierxsieben $  \\

\hline  \leitspaltefuenf  &  $ \afuenfxeins $  &  $ \afuenfxzwei $  &  $ \afuenfxdrei $  &  $ \afuenfxvier $ &  $ \afuenfxfuenf $  &  $ \afuenfxsechs $  &  $ \afuenfxsieben $  \\

\hline  \leitspaltesechs  &  $ \asechsxeins $  &  $ \asechsxzwei $  &  $ \asechsxdrei $  &  $ \asechsxvier $  &  $ \asechsxfuenf $  &  $ \asechsxsechs $  &  $ \asechsxsieben $  \\

\hline  \leitspaltesieben  &  $ \asiebenxeins $  &  $ \asiebenxzwei $  &  $ \asiebenxdrei $  &  $ \asiebenxvier $  &  $ \asiebenxfuenf $  &  $ \asiebenxsechs $  &  $ \asiebenxsieben $   \\

\hline  \leitspalteacht  &  $ \aachtxeins $  &  $ \aachtxzwei $  &  $ \aachtxdrei $  &  $ \aachtxvier $  &  $ \aachtxfuenf $  &  $ \aachtxsechs $  &  $ \aachtxsieben $ \\

\hline  \leitspalteneun  &  $ \aneunxeins $  &  $ \aneunxzwei $  &  $ \aneunxdrei $  &  $ \aneunxvier $  &  $ \aneunxfuenf $  &  $ \aneunxsechs $  &  $ \aneunxsieben $  \\

\hline  \leitspaltezehn  &  $ \azehnxeins $  &  $ \azehnxzwei $  &  $ \azehnxdrei $  &  $ \azehnxvier $ &  $ \azehnxfuenf $  &  $ \azehnxsechs $  &  $ \azehnxsieben $   \\

\hline  \leitspalteelf  &  $ \aelfxeins $  &  $ \aelfxzwei $  &  $ \aelfxdrei $  &  $ \aelfxvier $  &  $ \aelfxfuenf $  &  $ \aelfxsechs $  &  $ \aelfxsieben $  \\

\hline  \leitspaltezwoelf  &  $ \azwoelfxeins $  &  $ \azwoelfxzwei $  &  $ \azwoelfxdrei $  &  $ \zwoelfxvier $  &  $ \azwoelfxfuenf $  &  $ \azwoelfxsechs $  &  $ \zwoelfxsieben $ \\

\hline

\end{tabular}

\end{center} }

%Tabellen mit sechzehn Zeilen

\newcommand{\tabelleleitsechzehnxzwei}{

\begin{center}

\begin{tabular}{|c||c|c|}

\hline  \leitzeilenull  &  \leitzeileeins  &  \leitzeilezwei   \\

\hline \hline  \leitspalteeins  &  $ \aeinsxeins $  &  $ \aeinsxzwei $ \\

\hline  \leitspaltezwei  &  $ \azweixeins $  &  $ \azweixzwei $  \\

\hline  \leitspaltedrei  &  $ \adreixeins $  &  $ \adreixzwei $  \\

\hline  \leitspaltevier  &  $ \avierxeins $  &  $ \avierxzwei $  \\

\hline  \leitspaltefuenf  &  $ \afuenfxeins $  &  $ \afuenfxzwei $  \\

\hline  \leitspaltesechs  &  $ \asechsxeins $  &  $ \asechsxzwei $  \\

\hline  \leitspaltesieben  &  $ \asiebenxeins $  &  $ \asiebenxzwei $  \\

\hline  \leitspalteacht  &  $ \aachtxeins $  &  $ \aachtxzwei $  \\

\hline  \leitspalteneun  &  $ \aneunxeins $  &  $ \aneunxzwei $   \\

\hline  \leitspaltezehn  &  $ \azehnxeins $  &  $ \azehnxzwei $  \\

\hline  \leitspalteelf  &  $ \aelfxeins $  &  $ \aelfxzwei $   \\

\hline  \leitspaltezwoelf  &  $ \azwoelfxeins $  &  $ \azwoelfxzwei $   \\

\hline  \leitspaltedreizehn  &  $ \adreizehnxeins $  &  $ \adreizehnxzwei $   \\

\hline  \leitspaltevierzehn  &  $ \avierzehnxeins $  &  $ \avierzehnxzwei $  \\

\hline  \leitspaltefuenfzehn  &  $ \afuenfzehnxeins $  &  $ \afuenfzehnxzwei $   \\

\hline  \leitspaltesechzehn  &  $ \asechzehnxeins $  &  $ \asechzehnxzwei $   \\

\hline

\end{tabular}

\end{center} }

\newcommand{\tabelleleitsechzehnxsechs}{

\begin{center}

\begin{tabular}{|c||c|c|c|c|c|c|}

\hline  \leitzeilenull  &  \leitzeileeins  &  \leitzeilezwei  &  \leitzeiledrei  &  \leitzeilevier  &  \leitzeilefuenf  &  \leitzeilesechs  \\

\hline \hline  \leitspalteeins  &  $ \aeinsxeins $  &  $ \aeinsxzwei $  &  $ \aeinsxdrei $  &  $ \aeinsxvier $  &  $ \aeinsxfuenf $  &  $ \aeinsxsechs $  \\

\hline  \leitspaltezwei  &  $ \azweixeins $  &  $ \azweixzwei $  &  $ \azweixdrei $  &  $ \azweixvier $  &  $ \azweixfuenf $  &  $ \azweixsechs $  \\

\hline  \leitspaltedrei  &  $ \adreixeins $  &  $ \adreixzwei $  &  $ \adreixdrei $  &  $ \adreixvier $  &  $ \adreixfuenf $  &  $ \adreixsechs $  \\

\hline  \leitspaltevier  &  $ \avierxeins $  &  $ \avierxzwei $  &  $ \avierxdrei $  &  $ \avierxvier $  &  $ \avierxfuenf $  &  $ \avierxsechs $  \\

\hline  \leitspaltefuenf  &  $ \afuenfxeins $  &  $ \afuenfxzwei $  &  $ \afuenfxdrei $  &  $ \afuenfxvier $  &  $ \afuenfxfuenf $  &  $ \afuenfxsechs $  \\

\hline  \leitspaltesechs  &  $ \asechsxeins $  &  $ \asechsxzwei $  &  $ \asechsxdrei $  &  $ \asechsxvier $  &  $ \asechsxfuenf $  &  $ \asechsxsechs $  \\

\hline  \leitspaltesieben  &  $ \asiebenxeins $  &  $ \asiebenxzwei $  &  $ \asiebenxdrei $  &  $ \asiebenxvier $  &  $ \asiebenxfuenf $  &  $ \asiebenxsechs $  \\

\hline  \leitspalteacht  &  $ \aachtxeins $  &  $ \aachtxzwei $  &  $ \aachtxdrei $  &  $ \aachtxvier $  &  $ \aachtxfuenf $  &  $ \aachtxsechs $  \\

\hline  \leitspalteneun  &  $ \aneunxeins $  &  $ \aneunxzwei $  &  $ \aneunxdrei $  &  $ \aneunxvier $  &  $ \aneunxfuenf $  &  $ \aneunxsechs $  \\

\hline  \leitspaltezehn  &  $ \azehnxeins $  &  $ \azehnxzwei $  &  $ \azehnxdrei $  &  $ \azehnxvier $  &  $ \azehnxfuenf $  &  $ \azehnxsechs $  \\

\hline  \leitspalteelf  &  $ \aelfxeins $  &  $ \aelfxzwei $  &  $ \aelfxdrei $  &  $ \aelfxvier $  &  $ \aelfxfuenf $  &  $ \aelfxsechs $  \\

\hline  \leitspaltezwoelf  &  $ \azwoelfxeins $  &  $ \azwoelfxzwei $  &  $ \azwoelfxdrei $  &  $ \azwoelfxvier $  &  $ \azwoelfxfuenf $  &  $ \azwoelfxsechs $  \\

\hline  \leitspaltedreizehn  &  $ \adreizehnxeins $  &  $ \adreizehnxzwei $  &  $ \adreizehnxdrei $  &  $ \adreizehnxvier $  &  $ \adreizehnxfuenf $  &  $ \adreizehnxsechs $  \\

\hline  \leitspaltevierzehn  &  $ \avierzehnxeins $  &  $ \avierzehnxzwei $  &  $ \avierzehnxdrei $  &  $ \avierzehnxvier $  &  $ \avierzehnxfuenf $  &  $ \avierzehnxsechs $  \\

\hline  \leitspaltefuenfzehn  &  $ \afuenfzehnxeins $  &  $ \afuenfzehnxzwei $  &  $ \afuenfzehnxdrei $  &  $ \afuenfzehnxvier $  &  $ \afuenfzehnxfuenf $  &  $ \afuenfzehnxsechs $  \\

\hline  \leitspaltesechzehn  &  $ \asechzehnxeins $  &  $ \asechzehnxzwei $  &  $ \asechzehnxdrei $  &  $ \asechzehnxvier $  &  $ \asechzehnxfuenf $  &  $ \asechzehnxsechs $  \\

\hline

\end{tabular}

\end{center} }

%Tabellen mit achtzehn Zeilen

\newcommand{\tabelleleitachtzehnxneun}{

\begin{center}

\begin{tabular}{|c||c|c|c|c|c|c|c|c|c|}

\hline  \leitzeilenull  &  \leitzeileeins  &  \leitzeilezwei  &  \leitzeiledrei  &  \leitzeilevier  &  \leitzeilefuenf  &  \leitzeilesechs &  \leitzeilesieben  &  \leitzeileacht  &  \leitzeileneun \\

\hline \hline  \leitspalteeins  &  $ \aeinsxeins $  &  $ \aeinsxzwei $  &  $ \aeinsxdrei $  &  $ \aeinsxvier $  &  $ \aeinsxfuenf $  &  $ \aeinsxsechs $ &  $ \aeinsxsieben $  &  $ \aeinsxacht $  &  $ \aeinsxneun $  \\

\hline  \leitspaltezwei  &  $ \azweixeins $  &  $ \azweixzwei $  &  $ \azweixdrei $  &  $ \azweixvier $  &  $ \azweixfuenf $  &  $ \azweixsechs $ &  $ \azweixsieben $  &  $ \azweixacht $  &  $ \azweixneun $  \\

\hline  \leitspaltedrei  &  $ \adreixeins $  &  $ \adreixzwei $  &  $ \adreixdrei $  &  $ \adreixvier $  &  $ \adreixfuenf $  &  $ \adreixsechs $ &  $ \adreixsieben $  &  $ \adreixacht $  &  $ \adreixneun $  \\

\hline  \leitspaltevier  &  $ \avierxeins $  &  $ \avierxzwei $  &  $ \avierxdrei $  &  $ \avierxvier $  &  $ \avierxfuenf $  &  $ \avierxsechs $  &  $ \avierxsieben $  &  $ \avierxacht $  &  $ \avierxneun $ \\

\hline  \leitspaltefuenf  &  $ \afuenfxeins $  &  $ \afuenfxzwei $  &  $ \afuenfxdrei $  &  $ \afuenfxvier $  &  $ \afuenfxfuenf $  &  $ \afuenfxsechs $ &  $ \afuenfxsieben $  &  $ \afuenfxacht $  &  $ \afuenfxneun $  \\

\hline  \leitspaltesechs  &  $ \asechsxeins $  &  $ \asechsxzwei $  &  $ \asechsxdrei $  &  $ \asechsxvier $  &  $ \asechsxfuenf $  &  $ \asechsxsechs $ &  $ \asechsxsieben $  &  $ \asechsxacht $  &  $ \asechsxneun $  \\

\hline  \leitspaltesieben  &  $ \asiebenxeins $  &  $ \asiebenxzwei $  &  $ \asiebenxdrei $  &  $ \asiebenxvier $  &  $ \asiebenxfuenf $  &  $ \asiebenxsechs $ &  $ \asiebenxsieben $  &  $ \asiebenxacht $  &  $ \asiebenxneun $  \\

\hline  \leitspalteacht  &  $ \aachtxeins $  &  $ \aachtxzwei $  &  $ \aachtxdrei $  &  $ \aachtxvier $  &  $ \aachtxfuenf $  &  $ \aachtxsechs $ &  $ \aachtxsieben $  &  $ \aachtxacht $  &  $ \aachtxneun $  \\

\hline  \leitspalteneun  &  $ \aneunxeins $  &  $ \aneunxzwei $  &  $ \aneunxdrei $  &  $ \aneunxvier $  &  $ \aneunxfuenf $  &  $ \aneunxsechs $ &  $ \aneunxsieben $  &  $ \aneunxacht $  &  $ \aneunxneun $  \\

\hline  \leitspaltezehn  &  $ \azehnxeins $  &  $ \azehnxzwei $  &  $ \azehnxdrei $  &  $ \azehnxvier $  &  $ \azehnxfuenf $  &  $ \azehnxsechs $ &  $ \azehnxsieben $  &  $ \azehnxacht $  &  $ \azehnxneun $  \\

\hline  \leitspalteelf  &  $ \aelfxeins $  &  $ \aelfxzwei $  &  $ \aelfxdrei $  &  $ \aelfxvier $  &  $ \aelfxfuenf $  &  $ \aelfxsechs $ &  $ \aelfxsieben $  &  $ \aelfxacht $  &  $ \aelfxneun $  \\

\hline  \leitspaltezwoelf  &  $ \azwoelfxeins $  &  $ \azwoelfxzwei $  &  $ \azwoelfxdrei $  &  $ \azwoelfxvier $  &  $ \azwoelfxfuenf $  &  $ \azwoelfxsechs $ &  $ \azwoelfxsieben $  &  $ \azwoelfxacht $  &  $ \azwoelfxneun $  \\

\hline  \leitspaltedreizehn  &  $ \adreizehnxeins $  &  $ \adreizehnxzwei $  &  $ \adreizehnxdrei $  &  $ \adreizehnxvier $  &  $ \adreizehnxfuenf $  &  $ \adreizehnxsechs $ &  $ \adreizehnxsieben $  &  $ \adreizehnxacht $  &  $ \adreizehnxneun $  \\

\hline  \leitspaltevierzehn  &  $ \avierzehnxeins $  &  $ \avierzehnxzwei $  &  $ \avierzehnxdrei $  &  $ \avierzehnxvier $  &  $ \avierzehnxfuenf $  &  $ \avierzehnxsechs $  &  $ \avierzehnxsieben $  &  $ \avierzehnxacht $  &  $ \avierzehnxneun $ \\

\hline  \leitspaltefuenfzehn  &  $ \afuenfzehnxeins $  &  $ \afuenfzehnxzwei $  &  $ \afuenfzehnxdrei $  &  $ \afuenfzehnxvier $  &  $ \afuenfzehnxfuenf $  &  $ \afuenfzehnxsechs $ &  $ \afuenfzehnxsieben $  &  $ \afuenfzehnxacht $  &  $ \afuenfzehnxneun $  \\

\hline  \leitspaltesechzehn  &  $ \asechzehnxeins $  &  $ \asechzehnxzwei $  &  $ \asechzehnxdrei $  &  $ \asechzehnxvier $  &  $ \asechzehnxfuenf $  &  $ \asechzehnxsechs $ &  $ \asechzehnxsieben $  &  $ \asechzehnxacht $  &  $ \asechzehnxneun $  \\

\hline  \leitspaltesiebzehn  &  $ \asiebzehnxeins $  &  $ \asiebzehnxzwei $  &  $ \asiebzehnxdrei $  &  $ \asiebzehnxvier $  &  $ \asiebzehnxfuenf $  &  $ \asiebzehnxsechs $ &  $ \asiebzehnxsieben $  &  $ \asiebzehnxacht $  &  $ \asiebzehnxneun $  \\

\hline  \leitspalteachtzehn  &  $ \aachtzehnxeins $  &  $ \aachtzehnxzwei $  &  $ \aachtzehnxdrei $  &  $ \aachtzehnxvier $  &  $ \aachtzehnxfuenf $  &  $ \aachtzehnxsechs $ &  $ \aachtzehnxsieben $  &  $ \aachtzehnxacht $  &  $ \aachtzehnxneun $  \\

\hline

\end{tabular}

\end{center} }

\newcommand{\bruchhilfealign}{\cr & & }

\newcommand{\bruchhilfedisp}{\]  \[ }

\newcommand{\bruchhilfe}{ \- }

%[[Projekt:Semantische Vorlagen/Notlösungen/Latex]]

%
%Bei verschachtelten ifthen-Befehlen
%kann es Probleme geben, daher
%

\newcommand{\faktuebergangpos}[1]{ \faktuebergangskip  {\hspace{-0,55cm} } {#1} }

%Ende des Grundvorspannes

%Globale Strukturierungen

\newcommand{\seitenueberschrift}[1]{\seitenueberschriftvorskip
\begin{center} \large{ #1} \end{center} \seitenueberschriftnachskip}

\newcommand{\zwischenueberschrift}[1]{\section*{#1} }

\newcommand{\zwischenueberschriftaufgaben}[1]{ \zwischenueberschriftvorskip \begin{center} {\bf #1} \end{center} \zwischenueberschriftnachskip}

\newcommand{\unterueberschrift}[1]{ \unterueberschriftvorskip \hspace{0.5cm} {\em #1}  \unterueberschriftnachskip}

\renewcommand{\definitionbenennung}[1]{}

%Grund:Bei Definitionen im Normaltext
%nicht nötig, da Begriff früh kommt

\renewcommand{\beispielbenennung}[1]{}

\renewcommand{\bemerkungbenennung}[1]{}

%Umbrüche bei großer Schrift (entfällt)

%[[Projekt:Semantische Vorlagen/Mathematischer Modus/Umbrüche/Latex]]

\newcommand{\mathdisplaybruch}{}

\newcommand{\mathdisplaybruchinklammer}{}

\newcommand{\mathl}[2]{$#1$#2}

\newcommand{\mathlk}[2]{$#1$#2}

%Nicht durch semantische Vorlage produzierte Befehle

\newcommand{\mathdisplaybruchbruch}{ \]  \[ }

%Gestaltung der vertikalen und horizontalen Abstände

%[[Projekt:Semantische Vorlagen/Latex/Vorlesungsskript/Skips]]

%Abstände bei Überschriften

\newcommand{\seitenueberschriftvorskip}{\par \bigskip \bigskip \bigskip }

\newcommand{\seitenueberschriftnachskip}{\par \bigskip}

\newcommand{\zwischenueberschriftvorskip}{\par \bigskip}

\newcommand{\zwischenueberschriftnachskip}{\par \smallskip}

\newcommand{\unterueberschriftvorskip}{\par \bigskip}

\newcommand{\unterueberschriftnachskip}{\par \smallskip}

%Aufzählungen und Listen

\newcommand{\listskip}{\smallskip}

%Abstände bei mathematischen Strukturen

\newcommand{\aufgabevorskip}{\bigskip}

\newcommand{\aufgabepunktskip}{\par}

\newcommand{\aufgabenachskip}{}

\newcommand{\aufgabeloesungskip}{\par \bigskip}

\newcommand{\beispielvorskip}{}

\newcommand{\beispielnachskip}{}

\newcommand{\beispielbenennungnachskip}{}

\newcommand{\bemerkungvorskip}{}

\newcommand{\bemerkungnachskip}{}

\newcommand{\bemerkungbenennungnachskip}{}

\newcommand{\definitionvorskip}{}

\newcommand{\definitionnachskip}{}

\newcommand{\definitionbenennungnachskip}{}

\newcommand{\faktvorskip}{}

\newcommand{\faktnachskip}{}

\newcommand{\faktbenennungvorskip}{}

\newcommand{\faktbenennungnachskip}{}

\newcommand{\faktsituationskip}{}

\newcommand{\faktvoraussetzungskip}{}

\newcommand{\faktuebergangskip}{}

\newcommand{\faktfolgerungskip}{}

\newcommand{\faktzusatzskip}{}

%Sonstiges

\newcommand{\deckblattabstand}{\vspace{2cm}}

\newcommand{\gesamtskriptnewpage}{ }

\newcommand{\einzeltextnewpage}{ }

%[[Projekt:Semantische Vorlagen/Latex/Vorlesungsskript/Horizontale Abstände]]

\newcommand{\leerzeichen}{ }

%Seitengestaltung

%[[Projekt:Semantische Vorlagen/Latex/Vorlesungsskript/Seitengestaltung]]

%Einzelne Abstände

% page geometry supplied by reader wrapper
% page geometry supplied by reader wrapper

% page geometry supplied by reader wrapper
% page geometry supplied by reader wrapper

% page geometry supplied by reader wrapper
% page geometry supplied by reader wrapper

\setlength{\parindent}{0cm}

\setlength{\parskip}{1ex}

%Bildgestaltung

%[[Projekt:Semantische Vorlagen/Latex/Vorlesungsskript/Bildgestaltung]]

\newcommand{\bild}[1]{\vspace{4mm} #1}

\newcommand{\bildtext}[1]{\begin{center} {\small #1 } \end{center} }

%Lokale Einstellungen

%\input{Lokalername}

\newcommand{\bildeinlesung}[1]{../authority/media/#1}

\renewcommand{\bildeinlesung}[2]{../authority/media/#1.#2}

\newcommand{\bildeinlesungpng}[2]{../authority/media/#1.png}

\newcommand{\bildeinlesungPNG}[2]{../authority/media/#1.png}

\newcommand{\bildeinlesungjpg}[2]{../authority/media/#1.jpg}

\newcommand{\bildeinlesungJPG}[2]{../authority/media/#1.jpg}

\newcommand{\bildeinlesungjpeg}[2]{../authority/media/#1.jpg}

\newcommand{\bildeinlesungsvg}[2]{generated/media/#1.png}

\newcommand{\bildeinlesunggif}[2]{../authority/media/#1.png}

\newcommand{\bildeinlesungGIF}[2]{../authority/media/#1.png}

\newcommand{\bildeinlesungxcf}[2]{../authority/media/#1.png}

%Klausurgestaltung

%[[Projekt:Semantische Vorlagen/Klausurgestaltung/Latex]]

\newcommand{\fachbereich}{Fachbereich}

\newcommand{\dozent}{Dozent}

\newcommand{\klausurdatum}{Datum}

\newcommand{\klausurgebiet}{Mathematik}

\newcommand{\klausurtyp}{Klausur}

\newcommand{\klausurmodus}{Dauer: Zehn Minuten Einlesezeit und zwei volle Stunden Bearbeitungszeit. Es sind keine Hilfsmittel erlaubt. Alle Antworten sind zu begr\"unden.

Es gibt insgesamt 64 Punkte. Es gilt die Sockelregelung, d.h. die Bewertung pro Aufgabe(nteil) beginnt in der Regel bei der halben Punktzahl. Zum Bestehen braucht man $16$ Punkte, ab $32$ Punkten gibt es eine Eins.

Tragen Sie auf dem Deckblatt und jedem weiteren Blatt Ihren Namen und Ihre
Matrikelnummer leserlich ein. Viel Erfolg!}

\newcommand{\klausurname}{
\renewcommand{\arraystretch}{1.5}
\begin{tabular}{@{}p{3.2cm}p{10cm}}
Name, Vorname: &
..................................................................................\\
Matrikelnummer: &
..................................................................................\\
\end{tabular}
}

\newcommand{\klausurvorspann}[5]{\textbf{\fachbereich}\hfill \textbf{\klausurdatum}
\\
\textbf{\dozent} \hfill
\vspace{4ex}

\begin{center}
{\bfseries{\large \klausurgebiet \vspace{2ex}

\klausurtyp} }
\end{center}

\vspace{2ex}

\klausurmodus

\vspace{2ex}

\klausurname
}

\newcommand{\klausurnote}{\begin{tabular}{@{}p{3.2cm}p{10cm}} Note: & \end{tabular} }

\newcommand{\klausurtaeuschungwarnung}{ \vspace{2ex}{Mir ist bewu\ss t, dass bei einem schweren T\"auschungsversuch s\"amtliche Pr\"u\-fungsvorleis\-tungen in diesem Kurs verfallen. Ein schwerer T\"auschungsversuch liegt insbesondere beim Einsatz von Literatur und von elektronischen Hilfsmitteln vor. \vspace{1ex} } }

\newcommand{\klausureinverstaendnis}{
\vspace{2ex}
Ich erkl\"are mich durch meine Unterschrift einverstanden, dass mein Klausurergebnis mit meiner Matrikelnummer im Internet bekanntgegeben wird.
\vspace{1ex}

\begin{tabular}{@{}p{3.2cm}p{10cm}}
Unterschrift: &
.....................................................................................
\end{tabular}
}

%Variante für Testklausur

\newcommand{\testklausurvorspann}[5]{\textbf{\fachbereich}\hfill \textbf{\klausurdatum} \\ \textbf{\dozent} \hfill \vspace{4ex}

\begin{center} {\bfseries{\large \klausurgebiet \vspace{2ex}

\klausurtyp} } \end{center}

\vspace{2ex}

\testklausurmodus

\vspace{3ex}

\klausurname }

\newcommand{\testklausurmodus}{Dauer: Zehn Minuten Einlesezeit und zwei volle Stunden Bearbeitungszeit.

Es sind keine Hilfsmittel erlaubt.

Alle Antworten sind zu begr\"unden.

Es gibt insgesamt 64 Punkte. Es gilt die Sockelregelung, d.h. die Bewertung pro Aufgabe(nteil) beginnt in der Regel bei der halben Punktzahl.

Zur Orientierung: Zum Bestehen braucht man $16$ Punkte, ab $32$ Punkten gibt es eine Eins.

Tragen Sie auf dem Deckblatt und jedem weiteren Blatt Ihren Namen und Ihre Matrikelnummer leserlich ein.

Viel Erfolg!}

%\renewcommand{\inputaufgabegibtloesung}[4]{\aufgabevorskip \begin{aufgabe} \punkte {#1} \aufgabepunktskip #2 #3 #4\end{aufgabe} \aufgabenachskip}

\newcommand{\klausurmitloesungenvorspann}[5]{\textbf{\fachbereich}\hfill \textbf{\klausurdatum} \\ \textbf{\dozent} \hfill \vspace{4ex}

\begin{center} {\bfseries{\large \klausurgebiet \vspace{2ex}

\klausurtyp \hspace{0mm} mit L\"osungen} } \end{center}

}

%Lizenzierung und Bildverzeichnis

%[[Projekt:Semantische Vorlagen/Latex/Vorlesungsskript/Lizenzierung]]

\newcommand{\Lizenztext}{Teks sumber dibuat oleh Holger Brenner di Wikiversity berbahasa Jerman dan digunakan di sini berdasarkan CC BY-SA 4.0. Media mengikuti lisensi per berkas.}

\newcommand{\Abbildungslizenzierung}{ Erl\"auterung: Die in diesem Text verwendeten Bilder stammen aus Commons (also von http://commons.wikimedia.org) und haben eine Lizenz, die die Verwendung hier erlaubt. Die Bilder werden mit ihren Dateinamen auf Commons angef\"uhrt zusammen mit ihrem Autor bzw. Hochlader und der Lizenz. }

%Konfiguration des Bildverzeichnisses

\newcommand{\bildlizenz}[6]{
\ifthenelse {\equal{#2}{}}
{\addcontentsline{lof}{figure}{Sumber = #1, Kreator = Pengguna #3 di #4, Lisensi = #5 \bildlizenzskip  }}
{ \ifthenelse {\equal{#3}{}}
{ \addcontentsline{lof}{figure}{ Sumber = #1, Kreator = #2, Lisensi = #5 \bildlizenzskip }}
{ \addcontentsline{lof}{figure}{ Sumber = #1, Kreator = #2 (diunggah oleh Pengguna #3 di #4), Lisensi = #5 \bildlizenzskip }} } }

\newcommand{\bildlizenzskip}{ }

%Englische Variante

\newcommand{\ignoriereeng}[1]{}

\ignoriereeng{

%[[Projekt:Semantische Vorlagen/Latex/Englische Zusätze]]

\newtheorem{corollary}[fakt]{Corollary}

\newtheorem{Corollary}[fakt]{Corollary}

\theoremstyle{definition}

\newtheorem{example}[fakt]{Example}

\newtheorem{remark}[fakt]{Remark}

\newcommand{\inputremark}[2] {\bemerkungvorskip \begin{remark} \bemerkungbenennung{#1} \bemerkungbenennungnachskip #2 \end{remark} \bemerkungnachskip}

\newcommand{\inputexample}[2] {\beispielvorskip \begin{example} \beispielbenennung{#1} \beispielbenennungnachskip #2 \end{example} \beispielnachskip}

\newcommand{\inputfaktproof}[5] {\faktvorskip \begin{#2}

%\label{#1} (Absetzen, sonst kann das folgende dahinter rutschen)

\faktbenennung{#3} \faktbenennungnachskip #4 \end{#2} \faktnachskip \beweis{#5}}

\renewcommand{\proofname}{\hspace{-0.64cm}{\it Proof}}

}

%Variante für Artikel

\newcommand{\ignoriereart}[1]{}

\ignoriereart{
\newcommand{\titelueberschrift}[1]{\title[#1] { #1} \maketitle }

\renewcommand{\seitenueberschrift}[1]{ \section {#1} }

\renewcommand{\zwischenueberschrift}[1]{ \subsection {#1} }
}

% Indonesian aliases used by the independent translated reader.
\newtheorem{Teorema}[fakt]{Teorema}
\renewcommand{\proofname}{Bukti}

% The frozen source numbers facts and exercises by lecture/worksheet section.
% Book chapters are reader scaffolding and must not enter source identifiers.
\renewcommand{\thefakt}{\arabic{section}.\arabic{fakt}}

\addto\captionsindonesian{%
  \renewcommand{\contentsname}{Daftar Isi}%
  \renewcommand{\figurename}{Gambar}%
  \renewcommand{\listfigurename}{Daftar Gambar}%
}
\hypersetup{pdflang={id-ID},bookmarksopen=true}
\emergencystretch=1em
\ifdefined\pdfinfoomitdate\pdfinfoomitdate=1\fi
\ifdefined\pdftrailerid\pdftrailerid{}\fi
\ifdefined\pdfsuppressptexinfo\pdfsuppressptexinfo=15\fi

\title{Geometri Diferensial dan Manifold Mulus\\\large Pembaca kumulatif hingga Unit 3}
\author{Holger Brenner\\Terjemahan Bahasa Indonesia independen}
\date{Sumber: \textit{Differentialgeometrie (Osnabrück 2023)}}

\begin{document}
\hypersetup{pageanchor=false}
\maketitle
\hypersetup{pageanchor=true}
\frontmatter
\chapter*{Tentang edisi ini}
Pembaca kumulatif ini menerjemahkan Kuliah 1--3 dan Lembar Kerja 1--3 dari kursus Holger Brenner di Wikiversity berbahasa Jerman. Teks sumber digunakan berdasarkan CC BY-SA 4.0. Terjemahan ini merupakan karya independen dan bukan edisi resmi atau dukungan dari penulis maupun Wikiversity. Setiap gambar mengikuti status hak atau lisensi berkasnya sendiri, yang dicatat pada daftar gambar dan manifest media.

Sumber permanen dan hash revisi dicatat dalam berkas kontrol edisi ini. Rumus, urutan, latihan, penanda solusi, dan atribusi media dipertahankan; ID mesin hanya merupakan lapisan tambahan. Bagian solusi hanya memuat solusi yang benar-benar disediakan oleh sumber.
\tableofcontents

\mainmatter
\part{Unit 1}
\chapter[Kuliah 1]{Kuliah 1: Analisis dan Geometri}
\input{generated/lecture01.id.build.tex}

\chapter{Lembar Kerja 1}
\setcounter{section}{1}

  \zwischenueberschrift{Soal latihan}

\inputaufgabegibtloesung
{}
{

Misalkan

\mavergleichskette
{\vergleichskette
{ G
}
{ \subseteq }{ \R^n
}
{  }{
}
{    }{
}
{   }{
}
}
{}{}{}
\definitionsverweis {terbuka}{}{}
dan
\maabbdisp {\varphi} { G } { \R^m
} {}
sebuah
\definitionsverweis {pemetaan terdiferensialkan secara kontinu}{}{,}
yang di titik

\mavergleichskette
{\vergleichskette
{ P
}
{ \in }{ G
}
{  }{
}
{    }{
}
{   }{
}
}
{}{}{}
mempunyai
\definitionsverweis {diferensial total}{}{}
yang surjektif. Misalkan

\mavergleichskette
{\vergleichskette
{ v
}
{ \in }{ T_PZ
}
{  }{
}
{    }{
}
{   }{
}
}
{}{}{}
sebuah vektor di
\definitionsverweis {ruang tangen pada serat}{}{}
$Z$ dari $\varphi$ yang melalui $P$. Buktikan bahwa terdapat
\definitionsverweis {kurva terdiferensialkan secara kontinu}{}{}
\maabbdisp {\gamma} { ]- \epsilon, \epsilon[ } { Z
} {}
\zusatzklammer {untuk suatu

\mavergleichskettek
{\vergleichskettek
{ \epsilon
}
{ > }{ 0
}
{  }{
}
{    }{
}
{   }{
}
}
{}{}{}
yang sesuai} {} {}
dengan

\mavergleichskette
{\vergleichskette
{ \gamma(0)
}
{ = }{ P
}
{  }{
}
{    }{
}
{   }{
}
}
{}{}{}
dan

\mavergleichskettedisp
{\vergleichskette
{ \gamma'(0)
}
{ =} { v
}
{ } {
}
{ } {
}
{ } {
}
}
{}{}{.}

}
{} {}

\inputaufgabe
{}
{

Misalkan

\mavergleichskette
{\vergleichskette
{ W
}
{ \subseteq }{ \R^n
}
{  }{
}
{    }{
}
{   }{
}
}
{}{}{}
\definitionsverweis {terbuka}{}{,}
\maabb {h} { W } { \R
} {}
sebuah
\definitionsverweis {fungsi terdiferensialkan secara kontinu}{}{}
dan

\mavergleichskette
{\vergleichskette
{ Y
}
{ = }{ h^{-1}(0)
}
{  }{
}
{    }{
}
{   }{
}
}
{}{}{}
merupakan
\definitionsverweis {serat}{}{}
di atas $0$, dengan $h$
\definitionsverweis {reguler}{}{}
di setiap titik pada $Y$. Misalkan
\maabbdisp {g} { W } {\R
} {}
sebuah fungsi terdiferensialkan secara kontinu yang tidak mempunyai titik nol pada $Y$. Buktikan bahwa $Y$ juga dapat diperoleh sebagai serat dari

\mavergleichskettedisp
{\vergleichskette
{ \tilde{h}
}
{ =} { g h
}
{ } {
}
{ } {
}
{ } {
}
}
{}{}{,}
bahwa
\definitionsverweis {gradien}{}{}
dari $h$ dan $\tilde{h}$ merupakan kelipatan skalar satu sama lain, dan bahwa
\definitionsverweis {ruang-ruang tangen}{}{}
hanya bergantung pada $Y$.

}
{} {}

\inputaufgabe
{}
{

Misalkan

\mavergleichskette
{\vergleichskette
{ W
}
{ \subseteq }{ \R^n
}
{  }{
}
{    }{
}
{   }{
}
}
{}{}{}
\definitionsverweis {terbuka}{}{,}
\maabb {h} { W } { \R
} {}
sebuah
\definitionsverweis {fungsi terdiferensialkan secara kontinu}{}{}
dan

\mavergleichskette
{\vergleichskette
{ Y
}
{ = }{ h^{-1}(c)
}
{  }{
}
{    }{
}
{   }{
}
}
{}{}{}
merupakan
\definitionsverweis {serat}{}{}
di atas

\mavergleichskette
{\vergleichskette
{ c
}
{ \in }{ \R
}
{  }{
}
{    }{
}
{   }{
}
}
{}{}{,}
dengan $h$
\definitionsverweis {reguler}{}{}
di setiap titik pada $Y$. Misalkan
\maabbdisp {L} {\R^n } { \R^n
} {}
sebuah
\definitionsverweis {isomorfisme linear}{}{,}

\mavergleichskette
{\vergleichskette
{ W'
}
{ = }{ L^{-1}(W)
}
{  }{
}
{    }{
}
{   }{
}
}
{}{}{}
dan

\mavergleichskettedisp
{\vergleichskette
{ Z
}
{ =} { L^{-1}(Y)
}
{ =} { (h \circ L)^{-1}(c)
}
{ } {
}
{ } {
}
}
{}{}{.}
Buktikan bahwa konsep
\definitionsverweis {ruang tangen}{}{,}
medan vektor tangensial, dan solusi persamaan diferensial tangensial yang terkait pada $Y$ dan $Z$ secara umum saling bersesuaian
\zusatzklammer {yakni dapat dipetakan satu sama lain dengan bantuan $L$} {} {.}
Bagaimana dengan gradiennya?

}
{} {}

\inputaufgabe
{}
{

Misalkan

\mavergleichskette
{\vergleichskette
{ Y
}
{ = }{ h^{-1}(c)
}
{  }{
}
{    }{
}
{   }{
}
}
{}{}{}
sebuah
\definitionsverweis {hipermuka terdiferensialkan}{}{}
dan misalkan
\maabbdisp {f} { U} { \R
} {}
sebuah fungsi terdiferensialkan secara kontinu yang didefinisikan pada suatu lingkungan terbuka

\mavergleichskette
{\vergleichskette
{ Y
}
{ \subseteq }{ U
}
{  }{
}
{    }{
}
{   }{
}
}
{}{}{}
Andaikan $f$ di

\mavergleichskette
{\vergleichskette
{ P
}
{ \in }{ Y
}
{  }{
}
{    }{
}
{   }{
}
}
{}{}{}
mempunyai
\definitionsverweis {ekstremum lokal}{}{}
pada $Y$. Buktikan bahwa
\definitionsverweis {komponen tangensial}{}{}
dari $\operatorname{Grad} \,  f  (  P )$ sama dengan $0$.

}
{} {}

\inputaufgabe
{}
{

Misalkan

\mavergleichskette
{\vergleichskette
{ Y
}
{ \subseteq }{ \R^3
}
{  }{
}
{    }{
}
{   }{
}
}
{}{}{}
adalah permukaan bola satuan. Tentukan
\definitionsverweis {dekomposisi ortogonal}{}{}
vektor

\mavergleichskette
{\vergleichskette
{ \begin{pmatrix}  6  \\ 7 \\  -3   \end{pmatrix}
}
{ \in }{ \R^3
}
{  }{
}
{    }{
}
{   }{
}
}
{}{}{}
menjadi komponen tangensial dan komponen ortogonal di titik

\mavergleichskette
{\vergleichskette
{ \begin{pmatrix}  -  { \frac{   1  }{  2 } }  \\ { \frac{   1  }{  2 } } \\  { \frac{   1  }{  \sqrt{2}  } }   \end{pmatrix}
}
{ \in }{ Y
}
{  }{
}
{    }{
}
{   }{
}
}
{}{}{.}

}
{} {}

\inputaufgabe
{}
{

Misalkan

\mavergleichskette
{\vergleichskette
{ Y
}
{ \subseteq }{ \R^3
}
{  }{
}
{    }{
}
{   }{
}
}
{}{}{}
adalah
\definitionsverweis {permukaan bola satuan}{}{.}
Misalkan
\maabbdisp {\gamma} {[0, 2 \pi]} { Y
} {}
merupakan gerak melingkar beraturan yang mengelilingi permukaan bola pada ketinggian ${ \frac{   1  }{  2 } }$.
\aufzaehlungzwei {Parametrisasikan $\gamma$.
} {Tentukan
\definitionsverweis {percepatan tangensial}{}{}
dari $\gamma$ pada setiap saat

\mavergleichskette
{\vergleichskette
{ t
}
{ \in }{ [0, 2 \pi]
}
{  }{
}
{    }{
}
{   }{
}
}
{}{}{.}
}

}
{} {}

\inputaufgabe
{}
{

Misalkan $Y$ adalah
\definitionsverweis {grafik}{}{}
dari
\maabbeledisp {f} {\R^2} { \R
} {(x,y)} { x^2+y^3
} {,}
dan misalkan
\maabbdisp {\gamma} {\R} { Y
} {}
adalah lintasan pada grafik tersebut yang berada di atas kurva dasar
\mathl{t \mapsto (t,t)}{.} Tentukan, untuk setiap

\mavergleichskette
{\vergleichskette
{ t
}
{ \in }{ \R
}
{  }{
}
{    }{
}
{   }{
}
}
{}{}{,}
\definitionsverweis {percepatan tangensial}{}{}
dari $\gamma$ di $\gamma(t)$.

}
{} {}

\inputaufgabe
{}
{

Misalkan

\mavergleichskette
{\vergleichskette
{ Y
}
{ \subseteq }{ \R^n
}
{  }{
}
{    }{
}
{   }{
}
}
{}{}{}
sebuah hiperbidang, yaitu serat dari fungsi linear

\mavergleichskettedisp
{\vergleichskette
{ h(x_1 , \ldots , x_n)
}
{ =} { a_1x_1 + \cdots + a_nx_n
}
{ } {
}
{ } {
}
{ } {
}
}
{}{}{}
di atas

\mavergleichskette
{\vergleichskette
{ c
}
{ \in }{ \R
}
{  }{
}
{    }{
}
{   }{
}
}
{}{}{}
dengan

\mavergleichskette
{\vergleichskette
{ a_i
}
{ \neq }{ 0
}
{  }{
}
{    }{
}
{   }{
}
}
{}{}{}
untuk setidaknya satu indeks $i$. Tafsirkan pengertian
\definitionsverweis {diferensial total}{}{,}
\definitionsverweis {gradien}{}{,}
\definitionsverweis {ruang tangen}{}{,}
dekomposisi menjadi komponen tangensial dan normal,
\definitionsverweis {percepatan tangensial}{}{,}
serta
\definitionsverweis {kurva geodesik}{}{}
dalam kasus khusus ini.

}
{} {}

\inputaufgabe
{}
{

Misalkan

\mavergleichskette
{\vergleichskette
{ Y
}
{ = }{ h^{-1}(c)
}
{  }{
}
{    }{
}
{   }{
}
}
{}{}{}
sebuah
\definitionsverweis {hipermuka terdiferensialkan}{}{,}
yang memuat garis
\mathl{P+ \R v}{.} Buktikan bahwa garis yang dilalui secara seragam tersebut merupakan
\definitionsverweis {kurva geodesik}{}{}
pada $Y$.

}
{} {}

Dua titik

\mathbed {P,Q \in K} {}
 {P \neq Q} {}
 {} {} {} {,}
disebut \stichwort {antipodal} {,} jika garis yang menghubungkannya melalui pusat bola.

\inputaufgabe
{}
{

Misalkan

\mavergleichskette
{\vergleichskette
{ K
}
{ \subseteq }{ \R^3
}
{  }{
}
{    }{
}
{   }{
}
}
{}{}{}
sebuah permukaan bola. Buktikan bahwa setiap dua
\definitionsverweis {titik yang tidak antipodal}{}{}

\mathbed {P,Q \in K} {}
 {P \neq Q} {}
 {} {} {} {,}
terletak pada tepat satu
\definitionsverweis {lingkaran besar}{}{}
dari $K$.

}
{} {}

\bild{ \begin{center}
\includegraphics[width=5.5cm]{ \bildeinlesungsvg {Planned flight map of the Oiseau Blanc} {svg} }
\end{center}
\bildtext {Mengapa pesawat terbang menempuh lintasan melengkung?} }

\bildlizenz { Planned flight map of the Oiseau Blanc.svg } {} {Pethrus} {Commons} {CC BY-SA 3.0} {}

\inputaufgabe
{}
{

Sebuah pesawat akan terbang dari Osnabrück menuju suatu tempat tujuan di belahan Bumi selatan. Mungkinkah rutenya lebih pendek jika pesawat itu mula-mula terbang ke arah utara?

}
{} {}

\inputaufgabe
{}
{

Lucy Sonnenschein merasa bahwa hari-hari terlalu singkat. Karena itu, setiap hari dan selama masih terang, ia memutuskan terbang ke arah barat melintasi satu zona waktu agar setiap harinya berlangsung $25$, bukan $24$, jam.
\aufzaehlungfuenf{Dapatkah Lucy Sonnenschein meningkatkan proporsi jam terang dalam hidupnya dengan strategi ini?
}{Dapatkah Lucy Sonnenschein memperpanjang hidupnya dengan strategi ini?
}{Sekarang Lucy mempunyai teman di seluruh dunia. Setelah berapa hari ia bertemu dengan mereka lagi?
}{Apa yang menyebabkan Lucy mengalami kehilangan waktu yang mengimbangi tambahan waktu hariannya?
}{Apakah hal ini berkaitan dengan teori relativitas?
}

}
{} {}

\inputaufgabe
{}
{

Berikan contoh kurva dua kali terdiferensialkan yang bergerak pada suatu
\definitionsverweis {lingkaran besar}{}{}
dari permukaan bola satuan, tetapi bukan
\definitionsverweis {kurva geodesik}{}{.}

}
{} {}

\inputaufgabe
{}
{

Sebuah \stichwort {segitiga geodesik} {} pada permukaan bola

\mavergleichskettedisp
{\vergleichskette
{ S^2
}
{ =} { \left\{ (x,y,z) \in \R^3  \mid   x^2+y^2+z^2 = 1        \right\}
}
{ } {
}
{ } {
}
{ } {
}
}
{}{}{}
terdiri atas tiga titik berbeda yang tidak terletak pada satu lingkaran besar,

\mavergleichskette
{\vergleichskette
{ A,B,C
}
{ \in }{ S^2
}
{  }{
}
{    }{
}
{   }{
}
}
{}{}{,}
dengan setiap pasangan titik dihubungkan oleh busur yang lebih pendek dari suatu
\definitionsverweis {lingkaran besar}{}{}
\zusatzklammer {setiap pasangan titik tidak antipodal} {} {.}
Buktikan bahwa jumlah sudut dalam segitiga itu lebih besar daripada $\pi$.

}
{} {}

  \zwischenueberschrift{Soal untuk dikumpulkan}

\inputaufgabe
{3}
{

Misalkan $Y$ adalah hipermuka yang diberikan oleh syarat

\mavergleichskettedisp
{\vergleichskette
{ 2x^2 +3y^2+5z^2
}
{ =} { 10
}
{ } {
}
{ } {
}
{ } {
}
}
{}{}{}
di $\R^3$. Tentukan dekomposisi ortogonal vektor

\mavergleichskette
{\vergleichskette
{ \begin{pmatrix}  4  \\ 2 \\  -5   \end{pmatrix}
}
{ \in }{ \R^3
}
{  }{
}
{    }{
}
{   }{
}
}
{}{}{}
menjadi komponen tangensial dan komponen ortogonal di titik

\mavergleichskette
{\vergleichskette
{ \begin{pmatrix}  -1 \\ -1\\  1   \end{pmatrix}
}
{ \in }{ Y
}
{  }{
}
{    }{
}
{   }{
}
}
{}{}{.}

}
{} {}

\inputaufgabe
{4}
{

Misalkan $Y$ adalah
\definitionsverweis {grafik}{}{}
dari
\maabbeledisp {f} {\R^2} { \R
} {(x,y)} { x^2+y^2
} {,}
dan misalkan
\maabbdisp {\gamma} {\R} { Y
} {}
adalah lintasan pada grafik tersebut yang berada di atas kurva dasar
\mathl{x \mapsto (x,1)}{.} Tentukan, untuk setiap

\mavergleichskette
{\vergleichskette
{ x
}
{ \in }{ \R
}
{  }{
}
{    }{
}
{   }{
}
}
{}{}{,}
\definitionsverweis {percepatan tangensial}{}{}
dari $\gamma$ di $\gamma(x)$.

}
{} {}

\inputaufgabe
{4}
{

Pada 11 April 2023 pukul 17.45, secara tiba-tiba dan mengejutkan semua orang, rotasi Bumi pada porosnya dihentikan. Pada saat yang sama, hambatan udara dan semua kerugian akibat gesekan ditiadakan. Semua manusia, hewan, dan benda yang tidak terpasang tetap atau tidak sempat berpegangan akan meneruskan gerak sesaatnya sesuai hukum inersia. Gravitasi tetap bekerja sehingga, untungnya, mereka tidak terlepas ke angkasa. Kapan penghapus papan tulis dari ruang kuliah di Osnabrück
\zusatzklammer {penghapus itu terbang keluar melalui jendela} {} {}
untuk pertama kalinya kembali ke posisi awalnya?

}
{} {}

\inputaufgabe
{4 (3+1)}
{

Misalkan

\mavergleichskette
{\vergleichskette
{ W
}
{ \subseteq }{ \R^n
}
{  }{
}
{    }{
}
{   }{
}
}
{}{}{}
\definitionsverweis {terbuka}{}{,}
\maabb {h} { W } { \R
} {}
sebuah
\definitionsverweis {fungsi terdiferensialkan secara kontinu}{}{}
dan

\mavergleichskette
{\vergleichskette
{ Y
}
{ = }{ h^{-1}(c)
}
{  }{
}
{    }{
}
{   }{
}
}
{}{}{}
merupakan
\definitionsverweis {serat}{}{}
di atas

\mavergleichskette
{\vergleichskette
{ c
}
{ \in }{ \R
}
{  }{
}
{    }{
}
{   }{
}
}
{}{}{,}
dengan $h$
\definitionsverweis {reguler}{}{}
di setiap titik pada $Y$. Misalkan
\maabbdisp {\gamma} { I } { Y
} {}
sebuah kurva dua kali
\definitionsverweis {terdiferensialkan secara kontinu}{}{.}
\aufzaehlungzwei {Buktikan bahwa jika $\gamma$ merupakan
\definitionsverweis {kurva geodesik}{}{,}
maka
\definitionsverweis {norma}{}{}
dari $\gamma'$ adalah konstan.
} {Berikan contoh kurva dengan norma kecepatan konstan yang bukan kurva geodesik.
}

}
{} {}

\inputaufgabe
{3}
{

Misalkan
\maabb {h,g} {\R^n } { \R
} {}
adalah
\definitionsverweis {fungsi-fungsi terdiferensialkan secara kontinu}{}{}
dan

\mavergleichskette
{\vergleichskette
{ Y
}
{ = }{ h^{-1}(c)
}
{  }{
}
{    }{
}
{   }{
}
}
{}{}{}
serta

\mavergleichskette
{\vergleichskette
{ Z
}
{ = }{ g^{-1}(d)
}
{  }{
}
{    }{
}
{   }{
}
}
{}{}{}
adalah
\definitionsverweis {hipermuka terdiferensialkan}{}{}
yang terkait, dengan $h$ di setiap titik pada $Y$ dan $g$ di setiap titik pada $Z$ bersifat
\definitionsverweis {reguler}{}{.}
Misalkan

\mavergleichskette
{\vergleichskette
{ T
}
{ = }{ Y \cap Z
}
{  }{
}
{    }{
}
{   }{
}
}
{}{}{}
dan misalkan
\maabbdisp {\gamma} { I } { T
} {}
sebuah kurva dua kali terdiferensialkan secara kontinu yang, jika dipandang sebagai kurva pada $Y$, merupakan
\definitionsverweis {kurva geodesik}{}{.}
Apakah $\gamma$ juga merupakan kurva geodesik pada $Z$?

}
{} {}


\section*{Solusi yang disediakan oleh sumber}
\addcontentsline{toc}{section}{Solusi yang disediakan oleh sumber}
% Authority: German Wikiversity pageid 131353, revid 1111802, 2026-08-15T16:41:12Z.
% Exact UTF-8 witness SHA-256: b07aac01b6f803481fab9b5331e19480fb9bf0058597159b6152c65c7bc955ce.

\zwischenueberschrift{Solusi untuk Soal 1}

Berlaku

\mavergleichskettedisp
{\vergleichskette
{ v
}
{ \in }{ T_PZ
}
{ = }{ \ker \left((D\varphi)_P\right)
}
{ } {
}
{ } {
}
}
{}{}{.}
Menurut
\definitionsverweis {teorema fungsi implisit}{}{,}
terdapat himpunan terbuka

\mathbed {P \in W} {}
 {W \subseteq G} {}
 {} {} {} {,}
himpunan terbuka

\mavergleichskette
{\vergleichskette
{ V
}
{ \subseteq }{ \R^{n-m}
}
{ } {
}
{ } {
}
{ } {
}
}
{}{}{,}
dan pemetaan terdiferensialkan secara kontinu
\maabbdisp {\psi} { V } { W } {}
sedemikian sehingga

\mavergleichskette
{\vergleichskette
{ \psi(V)
}
{ \subseteq }{ Z \cap W
}
{ } {
}
{ } {
}
{ } {
}
}
{}{}{}
dan $\psi$ menginduksi
\definitionsverweis {bijeksi}{}{}
\maabbdisp {\psi} { V } { Z \cap W } {}.
Misalkan

\mavergleichskette
{\vergleichskette
{ Q
}
{ \in }{ V
}
{ } {
}
{ } {
}
{ } {
}
}
{}{}{}
adalah titik yang memenuhi

\mavergleichskettedisp
{\vergleichskette
{ \psi(Q)
}
{ =} { P
}
{ } {
}
{ } {
}
{ } {
}
}
{}{}{.}
Pemetaan $\psi$ di setiap titik

\mavergleichskette
{\vergleichskette
{ Q
}
{ \in }{ V
}
{ } {
}
{ } {
}
{ } {
}
}
{}{}{}
bersifat
\definitionsverweis {reguler}{}{}
dan
\definitionsverweis {diferensial total}{}{}
dari $\psi$ memenuhi

\mavergleichskettedisp
{\vergleichskette
{ (D\varphi)_P \circ (D\psi)_Q
}
{ =} { 0
}
{ } {
}
{ } {
}
{ } {
}
}
{}{}{.}
Dengan demikian,

\mavergleichskettedisp
{\vergleichskette
{ \operatorname{im} \left((D\psi)_Q\right)
}
{ =} { \ker \left((D\varphi)_P\right)
}
{ =} { T_PZ
}
{ } {
}
{ } {
}
}
{}{}{.}
Karena $\psi$ reguler di $Q$, pemetaan
\maabbdisp {(D\psi)_Q} { \R^{n-m} } { \R^n } {}
bersifat injektif, sedangkan pemetaan
\maabbdisp {(D\psi)_Q} { \R^{n-m} } { T_PZ } {}
bersifat bijektif. Misalkan

\mavergleichskette
{\vergleichskette
{ u
}
{ \in }{ \R^{n-m}
}
{ } {
}
{ } {
}
{ } {
}
}
{}{}{}
adalah prapeta dari $v$, dan misalkan
\maabbeledisp {\theta} { ]-\epsilon,\epsilon[ } { \R^{n-m} } {t} {Q+tu} {,}
dengan $\epsilon$ dipilih cukup kecil agar seluruh citra terletak di dalam $V$. Maka
\maabbdisp {\gamma=\psi\circ\theta} { ]-\epsilon,\epsilon[ } { \R^n } {}
mempunyai sifat

\mavergleichskettedisp
{\vergleichskette
{ \gamma(0)
}
{ =} { \psi(\theta(0))
}
{ =} { \psi(Q)
}
{ =} { P
}
{ } {
}
}
{}{}{}
dan

\mavergleichskettedisp
{\vergleichskette
{ \gamma'(0)
}
{ =} { (D\psi)_Q(\theta'(0))
}
{ =} { (D\psi)_Q(u)
}
{ =} { v
}
{ } {
}
}
{}{}{.}


\part{Unit 2}
\chapter[Kuliah 2]{Kuliah 2: Permukaan Putar, Medan Normal, dan Pemetaan Gauss}
\input{generated/lecture02.id.build.tex}

\chapter{Lembar Kerja 2}
\setcounter{section}{2}

  \zwischenueberschrift{Soal latihan}

\inputaufgabegibtloesung
{}
{

Misalkan
\maabbdisp {f,g} {\R} { \R^n
} {}
\definitionsverweis {kurva-kurva terdiferensialkan}{}{.}
Hitung
\definitionsverweis {turunan}{}{}
fungsi
\maabbeledisp {} {\R} {\R
} { t } { \left\langle f(t) , g(t)  \right\rangle
} {.}
Nyatakan hasilnya menggunakan hasil kali dalam.

}
{} {}

\inputaufgabegibtloesung
{}
{

Misalkan
\maabb {h} { I } { V
} {}
sebuah kurva yang terdiferensialkan dua kali dalam
\definitionsverweis {ruang vektor Euklides}{}{}
$V$. Buktikan bahwa jika

\mavergleichskette
{\vergleichskette
{h'(t)
}
{ \neq }{ 0
}
{  }{
}
{    }{
}
{   }{
}
}
{}{}{,}
maka berlaku persamaan

\mavergleichskettedisp
{\vergleichskette
{ \Vert { h'(t)} \Vert'
}
{ =} { { \frac{   \left\langle h'(t) , h^{\prime \prime} (t)  \right\rangle  }{  \left\langle h'(t) , h'(t)  \right\rangle  } } \Vert { h'(t)} \Vert
}
{ } {
}
{ } {
}
{ } {
}
}
{}{}{.}

}
{} {}

Sebuah
\definitionsverweis {kurva terdiferensialkan}{}{}
\maabbeledisp {f} { [a,b] } {\R^n
} { t } {f(t) = \left( f_1(t)  , \,   \ldots  , \,  f_n(t)                 \right)
} {,}
disebut
\definitionswort {berparametrisasi panjang busur}{,}
jika

\mavergleichskette
{\vergleichskette
{ f_1'(t)^2 + \cdots +  f_n'(t)^2
}
{ = }{ 1
}
{  }{
}
{    }{
}
{   }{
}
}
{}{}{}
untuk setiap $t$.

\inputaufgabe
{}
{

Misalkan
\maabb {\gamma} { I } {\R^n
} {}
sebuah kurva yang dua kali
\definitionsverweis {terdiferensialkan secara kontinu}{}{}
dan
\definitionsverweis {berparametrisasi panjang busur}{}{.}
Buktikan bahwa
\mathkor {} {\gamma'(t)} {dan} {\gamma^{\prime \prime}(t)} {}
saling tegak lurus.

}
{} {}

\inputaufgabe
{}
{

Misalkan

\mavergleichskette
{\vergleichskette
{ P
}
{ \in }{ \R^n
}
{  }{
}
{    }{
}
{   }{
}
}
{}{}{}
sebuah titik dan misalkan

\mavergleichskette
{\vergleichskette
{ I
}
{ = }{ {]{-1},1[}
}
{  }{
}
{    }{
}
{   }{
}
}
{}{}{.}
Kita tinjau himpunan

\mavergleichskettedisp
{\vergleichskette
{ M
}
{ =} { \left\{ f:I \rightarrow \R^n   \mid  f \text{ terdiferensialkan} , \,  f(0) = P       \right\}
}
{ } {
}
{ } {
}
{ } {
}
}
{}{}{.}
Dua
\definitionsverweis {kurva}{}{}

\mavergleichskette
{\vergleichskette
{ f,g
}
{ \in }{ M
}
{  }{
}
{    }{
}
{   }{
}
}
{}{}{}
disebut \stichwort {ekuivalen secara tangensial} {,} jika

\mavergleichskettedisp
{\vergleichskette
{ f'(0)
}
{ =} { g'(0)
}
{ } {
}
{ } {
}
{ } {
}
}
{}{}{.}

a) Buktikan bahwa ini merupakan suatu
\definitionsverweis {relasi ekuivalensi}{}{.}

b) Temukan wakil paling sederhana bagi setiap
\definitionsverweis {kelas ekuivalensi}{}{.}

c) Berikan satu wakil lain untuk setiap kelas.

d) Deskripsikan himpunan kelas-kelas ekuivalensi tersebut
\zusatzklammer {yakni
\definitionsverweis {himpunan hasil bagi}{}{}} {} {.}

}
{} {}

\inputaufgabe
{}
{

Dengan syarat
\zusatzklammer {reguler di setiap titik} {} {,}
berikan sebuah contoh
\definitionsverweis {hipermuka terdiferensialkan}{}{,}
yang tidak
\definitionsverweis {terhubung}{}{.}

}
{} {}

\inputaufgabe
{}
{

Misalkan

\mavergleichskette
{\vergleichskette
{ V
}
{ \subseteq }{ \R^{n-1}
}
{  }{
}
{    }{
}
{   }{
}
}
{}{}{}
sebuah
\definitionsverweis {himpunan terbuka}{}{}
dan misalkan
\maabbdisp {f} { V } {\R
} {}
sebuah fungsi yang
\definitionsverweis {terdiferensialkan secara kontinu}{}{.}
Kita tinjau
\definitionsverweis {grafik}{}{}
$Y$ dari $f$ sebagai hipermuka terdiferensialkan di $\R^n$ dalam pengertian
Contoh 1.2.
Tentukan
\definitionsverweis {medan-medan normal satuan}{}{}
dalam kasus ini.

}
{} {}

\inputaufgabegibtloesung
{}
{

Misalkan
\mathkor {} {r} {dan} {R} {}
bilangan real positif yang memenuhi

\mavergleichskette
{\vergleichskette
{ 0
}
{ < }{ r
}
{ < }{ R
}
{    }{
}
{   }{
}
}
{}{}{.}
Tentukan suatu
\definitionsverweis {medan normal satuan}{}{}
untuk torus
\zusatzklammer {tertanam} {} {}

\mavergleichskettedisp
{\vergleichskette
{ T
}
{ =} { \left\{ (x,y,z) \in \R^3  \mid   \left(  \sqrt{x^2+y^2} -R  \right)^2 +z^2  = r^2         \right\}
}
{ } {
}
{ } {
}
{ } {
}
}
{}{}{.}

}
{} {}

\inputaufgabe
{}
{

Misalkan
\maabbdisp {f} { I } {\R_+
} {}
sebuah
\definitionsverweis {fungsi terdiferensialkan}{}{}
dan misalkan $Y$ adalah permukaan putar yang terkait dengan grafik $f$, dipandang sebagai hipermuka terdiferensialkan di $\R^3$ dalam pengertian
Lemma 2.1.
Tentukan
\definitionsverweis {medan-medan normal satuan}{}{}
dalam kasus ini.

}
{} {}

\inputaufgabe
{}
{

Buktikan bahwa
\definitionsverweis {pemetaan Gauss}{}{}
pada sebuah hiperbidang bersifat konstan. Apakah pernyataan sebaliknya juga berlaku?

}
{} {}

\inputaufgabe
{}
{

Buktikan bahwa
\definitionsverweis {pemetaan Gauss}{}{}
pada sebuah permukaan bola di $\R^n$ yang berpusat di titik asal, untuk setiap pilihan orientasi, diinduksi oleh homoteti skalar pada $\R^n$ (dengan faktor positif untuk orientasi keluar dan faktor negatif untuk orientasi masuk).

}
{} {}

\inputaufgabe
{}
{

Misalkan

\mavergleichskette
{\vergleichskette
{ Y
}
{ \subseteq }{ \R^n
}
{  }{
}
{    }{
}
{   }{
}
}
{}{}{}
sebuah
\definitionsverweis {hipermuka terdiferensialkan}{}{}
dengan
\definitionsverweis {medan normal satuan}{}{}
$N$ dan medan normal satuan yang berlawanan, $-N$. Buktikan bahwa kedua
\definitionsverweis {pemetaan Gauss}{}{}
yang terkait saling berhubungan melalui
\definitionsverweis {pemetaan antipodal}{}{}
\maabbeledisp {} {S^{n-1}} {S^{n-1}
} { P } { -P
} {.}

}
{} {}

\inputaufgabegibtloesung
{}
{

Misalkan

\mavergleichskette
{\vergleichskette
{ \alpha,\beta
}
{ \in }{ \R_+
}
{  }{
}
{    }{
}
{   }{
}
}
{}{}{}
dan misalkan

\mavergleichskettedisp
{\vergleichskette
{ Y
}
{ =} { \left\{ (x,y)\in\R^2  \mid   \alpha x^2+\beta y^2 = 1        \right\}
}
{ } {
}
{ } {
}
{ } {
}
}
{}{}{.}
Orientasikan $Y$ dengan medan normal satuan yang diperoleh dari gradien keluar fungsi
$ (x,y) \longmapsto \alpha x^2+\beta y^2 $.
\aufzaehlungzwei {Tentukan
\definitionsverweis {pemetaan Gauss}{}{}
untuk $Y$.
} { Berikan secara eksplisit suatu
\definitionsverweis {pemetaan invers}{}{}
untuk pemetaan Gauss pada $Y$.
}

}
{} {}

\inputaufgabegibtloesung
{}
{

Misalkan

\mavergleichskette
{\vergleichskette
{ \alpha,\beta
}
{ \in }{ \R_+
}
{  }{
}
{    }{
}
{   }{
}
}
{}{}{}
dan misalkan

\mavergleichskettedisp
{\vergleichskette
{ Y
}
{ =} { \left\{ (x,y)\in\R^2  \mid   \alpha x^2+\beta y^2 = 1        \right\}
}
{ } {
}
{ } {
}
{ } {
}
}
{}{}{.}
Buktikan bahwa
\definitionsverweis {pemetaan Gauss}{}{}
untuk $Y$ bersifat bijektif.

}
{} {}

\inputaufgabe
{}
{

Berikan sebuah contoh
\definitionsverweis {hipermuka terdiferensialkan}{}{}

\mavergleichskette
{\vergleichskette
{ Y
}
{ \subseteq }{ \R^2
}
{  }{
}
{    }{
}
{   }{
}
}
{}{}{}
sedemikian sehingga
\definitionsverweis {pemetaan Gauss}{}{}
bersifat surjektif dan terdapat titik-titik pada $S^1$ yang dicapai tak berhingga kali.

}
{} {}

\inputaufgabe
{}
{

Misalkan

\mavergleichskette
{\vergleichskette
{ W
}
{ \subseteq }{ \R^n
}
{  }{
}
{    }{
}
{   }{
}
}
{}{}{}
\definitionsverweis {terbuka}{}{,}
\maabb {h} { W } { \R
} {}
sebuah
\definitionsverweis {fungsi yang terdiferensialkan secara kontinu}{}{}
dan

\mavergleichskette
{\vergleichskette
{ Y
}
{ = }{ h^{-1}(c)
}
{  }{
}
{    }{
}
{   }{
}
}
{}{}{}
adalah
\definitionsverweis {serat}{}{}
di atas

\mavergleichskette
{\vergleichskette
{ c
}
{ \in }{ \R
}
{  }{
}
{    }{
}
{   }{
}
}
{}{}{,}
dengan $h$
\definitionsverweis {reguler}{}{}
di setiap titik pada $Y$. Misalkan
\maabbdisp {L} {\R^n } { \R^n
} {}
sebuah
\definitionsverweis {isometri linear}{}{,}

\mavergleichskette
{\vergleichskette
{ W'
}
{ = }{ L^{-1}(W)
}
{  }{
}
{    }{
}
{   }{
}
}
{}{}{}
dan

\mavergleichskettedisp
{\vergleichskette
{ Z
}
{ =} { L^{-1}(Y)
}
{ =} { (h \circ L)^{-1}(c)
}
{ } {
}
{ } {
}
}
{}{}{.}
Buktikan bahwa konsep-konsep kurva geodesik, medan normal, medan normal satuan, dan pemetaan Gauss pada $Y$ dan $Z$ saling bersesuaian
\zusatzklammer {yakni dapat dipetakan satu sama lain dengan bantuan $L$} {} {.}

}
{} {}

  \zwischenueberschrift{Soal untuk dikumpulkan}

\inputaufgabe
{3}
{

Tentukan citra parabola standar berorientasi ke atas di bawah
\definitionsverweis {pemetaan Gauss}{}{.}

}
{} {}

\inputaufgabe
{4}
{

Misalkan

\mavergleichskette
{\vergleichskette
{ W
}
{ = }{ \R^3 \setminus \{ (0,0,0) \}
}
{  }{
}
{    }{
}
{   }{
}
}
{}{}{}
dan misalkan

\mavergleichskette
{\vergleichskette
{ Y
}
{ \subseteq }{ W
}
{  }{
}
{    }{
}
{   }{
}
}
{}{}{}
adalah
\definitionsverweis {himpunan nol}{}{}
dari

\mavergleichskettedisp
{\vergleichskette
{ h(x,y,z)
}
{ =} { x^2+y^2-z^2
}
{ } {
}
{ } {
}
{ } {
}
}
{}{}{.}
Tentukan citra $Y$
\zusatzklammer {dengan orientasi yang ditentukan oleh gradien $h$} {} {}
di bawah
\definitionsverweis {pemetaan Gauss}{}{.}

}
{} {}

\inputaufgabe
{5}
{

Misalkan

\mavergleichskette
{\vergleichskette
{ \alpha,\beta, \gamma
}
{ \in }{ \R_+
}
{  }{
}
{    }{
}
{   }{
}
}
{}{}{}
dan misalkan

\mavergleichskettedisp
{\vergleichskette
{ Y
}
{ =} { \left\{ (x,y,z) \in \R^3  \mid   \alpha x^2+ \beta y^2 + \gamma z^2 = 1        \right\}
}
{ } {
}
{ } {
}
{ } {
}
}
{}{}{.}
Buktikan bahwa
\definitionsverweis {pemetaan Gauss}{}{}
untuk $Y$ bersifat bijektif.

}
{} {}

\inputaufgabe
{3}
{

Misalkan

\mavergleichskette
{\vergleichskette
{ W
}
{ \subseteq }{ \R^n
}
{  }{
}
{    }{
}
{   }{
}
}
{}{}{}
\definitionsverweis {terbuka}{}{,}
\maabb {h} { W } { \R
} {}
sebuah
\definitionsverweis {fungsi yang terdiferensialkan secara kontinu}{}{}
dan

\mavergleichskette
{\vergleichskette
{ Y
}
{ = }{ h^{-1}(c)
}
{  }{
}
{    }{
}
{   }{
}
}
{}{}{}
adalah
\definitionsverweis {serat}{}{}
di atas

\mavergleichskette
{\vergleichskette
{ c
}
{ \in }{ \R
}
{  }{
}
{    }{
}
{   }{
}
}
{}{}{,}
dengan $h$
\definitionsverweis {reguler}{}{}
di setiap titik pada $Y$. Misalkan
\maabbdisp {L} {\R^n } { \R^n
} {}
sebuah
\definitionsverweis {isomorfisme linear}{}{,}

\mavergleichskette
{\vergleichskette
{ W'
}
{ = }{ L^{-1}(W)
}
{  }{
}
{    }{
}
{   }{
}
}
{}{}{}
dan

\mavergleichskettedisp
{\vergleichskette
{ Z
}
{ =} { L^{-1}(Y)
}
{ =} { (h \circ L)^{-1}(c)
}
{ } {
}
{ } {
}
}
{}{}{.}
Buktikan bahwa konsep-konsep kurva geodesik, medan normal, medan normal satuan, dan pemetaan Gauss pada $Y$ dan $Z$ secara umum tidak saling bersesuaian
\zusatzklammer {berbeda dengan
Soal 2.14,
yang melibatkan suatu isometri} {} {.}

}
{} {}


\section*{Solusi yang disediakan oleh sumber}
\addcontentsline{toc}{section}{Solusi yang disediakan oleh sumber}
\subsection*{Solusi Soal 2.1}
\input{generated/worksheet02_exercise01_solution.id.build.tex}
\subsection*{Solusi Soal 2.2}
% Authority: German Wikiversity pageid 120435, revid 1089268, 2026-05-31T10:02:12Z.
% Exact UTF-8 wikitext witness SHA-256: 0bdd1cde6d1d0dc74757aa824634d10b6b4bf413efb4431538b7d9bac2b6f06c.
% Expanded German TeX source SHA-256: a92e556f50d2216192fc61703eea4bae3d233ab632c8eedb19bd35dec2ed89b0.

\textit{Catatan edisi: solusi sumber memakai basis ortogonal seolah-olah
ortonormal dan menuliskan turunan koordinat yang keliru. Argumen berikut
mengganti langkah tersebut.}

Karena $h'(t)\neq 0$, norma kecepatannya tak nol. Turunkan kuadrat normanya:

\[
\frac{d}{dt}\Vert h'(t)\Vert^2
=2\Vert h'(t)\Vert\bigl(\Vert h'(t)\Vert\bigr)'
=2\left\langle h'(t),h''(t)\right\rangle .
\]

Membagi dengan $2\Vert h'(t)\Vert$ memberikan

\[
\bigl(\Vert h'(t)\Vert\bigr)'
=\frac{\left\langle h'(t),h''(t)\right\rangle}{\Vert h'(t)\Vert}
=\frac{\left\langle h'(t),h''(t)\right\rangle}
       {\left\langle h'(t),h'(t)\right\rangle}\Vert h'(t)\Vert .
\]

\subsection*{Solusi Soal 2.7}
\begingroup\scriptsize
\input{generated/worksheet02_exercise07_solution.id.build.tex}
\endgroup
\subsection*{Solusi Soal 2.12}
\begingroup\small
\input{generated/worksheet02_exercise12_solution.id.build.tex}
\endgroup
\subsection*{Solusi Soal 2.13}
\input{generated/worksheet02_exercise13_solution.id.build.tex}

\part{Unit 3}
\chapter[Kuliah 3]{Kuliah 3: Kelengkungan Kurva Berparameter Panjang Busur}
\input{generated/lecture03.id.build.tex}

\chapter{Lembar Kerja 3}
\setcounter{section}{3}

  \zwischenueberschrift{Soal latihan}

\inputaufgabe
{}
{

Misalkan
\maabb {\gamma} { I } {\R^2
} {}
sebuah
\definitionsverweis {kurva berparametrisasi panjang busur}{}{}
yang dua kali terdiferensialkan secara kontinu, dan misalkan
\maabbeledisp {\delta} { -I } {\R^2
} { s } { \delta(s) = \gamma( -s)
} {,}
yaitu kurva yang ditelusuri dalam arah sebaliknya. Buktikan bahwa
\definitionsverweis {kelengkungan}{}{}
$\delta$ di $s$ merupakan negatif dari kelengkungan
\mathl{\gamma}{} di $-s$.

}
{} {}

\inputaufgabe
{}
{

Misalkan
\maabb {\gamma} { I } {\R^2
} {}
sebuah
\definitionsverweis {kurva berparametrisasi panjang busur}{}{}
yang dua kali terdiferensialkan secara kontinu, dan
\maabbdisp {L} {\R^2} {\R^2
} {}
sebuah
\definitionsverweis {rotasi}{}{.}
Buktikan bahwa
\definitionsverweis {kelengkungan}{}{}
$\gamma$ sama dengan kelengkungan
\mathl{L\circ \gamma}{}.

}
{} {}

\inputaufgabe
{}
{

Buktikan
Lema 3.6
dalam kasus dengan orientasi nonstandar.

}
{} {}

\inputaufgabe
{}
{

Misalkan
\maabb {\gamma} { I } {\R^2
} {}
sebuah
\definitionsverweis {kurva berparametrisasi panjang busur}{}{}
yang dua kali terdiferensialkan secara kontinu, dan
\maabbdisp {L} {\R^2} {\R^2
} {}
sebuah
\definitionsverweis {homoteti skalar}{}{}
dengan faktor homoteti

\mavergleichskette
{\vergleichskette
{ s
}
{ \neq }{ 0
}
{  }{
}
{    }{
}
{   }{
}
}
{}{}{.}
Bagaimana hubungan antara
\definitionsverweis {kelengkungan}{}{}
$\gamma$ dan kelengkungan
\mathl{L\circ \gamma}{?}

}
{} {}

\inputaufgabe
{}
{

Misalkan
\maabbeledisp {\gamma} {\R} { \R^2
} { t } { \begin{pmatrix}  \cos \left(  \omega t \right)  \\ \sin  \left(  \omega t \right)    \end{pmatrix}
} {,}
dengan

\mavergleichskette
{\vergleichskette
{ \omega
}
{ > }{ 0
}
{  }{
}
{    }{
}
{   }{
}
}
{}{}{.}
Buktikan bahwa terdapat tak berhingga banyak gerak melingkar dengan norma kecepatan konstan yang memiliki nilai dan turunan pertama yang sama dengan $\gamma$ pada

\mavergleichskette
{\vergleichskette
{ t
}
{ = }{ 0
}
{  }{
}
{    }{
}
{   }{
}
}
{}{}{.}

}
{} {}

Soal berikut penting dalam pembuatan komidi putar. Kesenangan menaiki komidi putar berasal dari percepatan, bukan dari kecepatan.

\inputaufgabe
{}
{

Misalkan
\maabbeledisp {\gamma} {\R} { \R^2
} { t } { \begin{pmatrix}  \cos \left(  \omega t \right)  \\ \sin  \left(  \omega t \right)    \end{pmatrix}
} {,}
dengan

\mavergleichskette
{\vergleichskette
{ \omega
}
{ > }{ 0
}
{  }{
}
{    }{
}
{   }{
}
}
{}{}{.}
Buktikan bahwa terdapat tak berhingga banyak gerak melingkar dengan norma kecepatan konstan yang berimpit dengan $\gamma$ pada

\mavergleichskette
{\vergleichskette
{ t
}
{ = }{ 0
}
{  }{
}
{    }{
}
{   }{
}
}
{}{}{}
dan juga memiliki percepatan yang sama di sana.

}
{} {}

\inputaufgabegibtloesung
{}
{

Misalkan
\maabbeledisp {\gamma} {\R} { \R^2
} { t } { \begin{pmatrix}  \cos \left(  \omega t \right)  \\ \sin  \left(  \omega t \right)    \end{pmatrix}
} {,}
dengan

\mavergleichskette
{\vergleichskette
{ \omega
}
{ > }{ 0
}
{  }{
}
{    }{
}
{   }{
}
}
{}{}{.}
Buktikan bahwa tidak ada gerak melingkar lain dengan norma kecepatan konstan yang memiliki nilai serta turunan pertama dan kedua yang sama dengan $\gamma$ pada

\mavergleichskette
{\vergleichskette
{ t
}
{ = }{ 0
}
{  }{
}
{    }{
}
{   }{
}
}
{}{}{.}

}
{} {}

\inputaufgabe
{}
{

Misalkan

\mavergleichskette
{\vergleichskette
{ W
}
{ \subseteq }{ \R^2
}
{  }{
}
{    }{
}
{   }{
}
}
{}{}{}
\definitionsverweis {terbuka}{}{,}
\maabb {h} { W } { \R
} {}
sebuah fungsi yang dua kali
\definitionsverweis {terdiferensialkan secara kontinu}{}{,}
dan

\mavergleichskette
{\vergleichskette
{ Y
}
{ = }{ h^{-1}(c)
}
{  }{
}
{    }{
}
{   }{
}
}
{}{}{}
adalah
\definitionsverweis {serat}{}{}
di atas

\mavergleichskette
{\vergleichskette
{ c
}
{ \in }{ \R
}
{  }{
}
{    }{
}
{   }{
}
}
{}{}{,}
dengan $h$
\definitionsverweis {reguler}{}{}
di setiap titik pada $Y$. Buktikan bahwa
\maabbdisp {\gamma} { I } { Y
} {}
merupakan
\definitionsverweis {kurva geodesik}{}{}
jika dan hanya jika $\gamma$
\definitionsverweis {berparametrisasi panjang busur}{}{.}

}
{} {}

\inputaufgabe
{}
{

\bild{ \begin{center}
\includegraphics[width=5.5cm]{ \bildeinlesungsvg {Euler spiral} {svg} }
\end{center}
\bildtext {} }

\bildlizenz { Euler spiral.svg } {} {AdiJapan} {Commons} {CC-by-sa 3.0} {}

Kita tinjau kurva terdiferensialkan
\zusatzklammer {yang disebut \stichwort {klotoid} {}} {} {}
\maabbeledisp {} {\R} {\R^2
} { t } { \gamma(t)
} {,}
dengan

\mavergleichskettedisp
{\vergleichskette
{ \gamma(t)
}
{ =} { \sqrt{\pi} \int_0^t  \begin{pmatrix}  \cos  { \frac{   \pi s^2 }{  2 } }  \\ \sin  { \frac{   \pi s^2 }{  2 } }     \end{pmatrix}  ds
}
{ } {
}
{ } {
}
{ } {
}
}
{}{}{.}
Buktikan bahwa
\definitionsverweis {kelengkungan}{}{}
kurva ini memenuhi persamaan

\mavergleichskettedisp
{\vergleichskette
{ \kappa(t)
}
{ =} { \sqrt{\pi} t
}
{ } {
}
{ } {
}
{ } {
}
}
{}{}{.}

}
{} {}

\inputaufgabe
{}
{

Misalkan
\maabb {f} { I } {\R
} {}
sebuah fungsi yang dua kali
\definitionsverweis {terdiferensialkan secara kontinu}{}{.}
Tentukan
\definitionsverweis {kelengkungan}{}{}
dan lingkaran kelengkungan dari grafik yang terkait.

}
{} {}

\inputaufgabe
{}
{

Misalkan
\maabb {f} {\R} {\R
} {}
dua kali
\definitionsverweis {terdiferensialkan secara kontinu}{}{}
dan

\mavergleichskettedisp
{\vergleichskette
{ \alpha(t)
}
{ \defeq} { \begin{pmatrix}  \cos f(t)  \\ \sin f(t)   \end{pmatrix}
}
{ } {
}
{ } {
}
{ } {
}
}
{}{}{.}
Tentukan
\definitionsverweis {kelengkungan}{}{}
$\alpha$ dan
\definitionsverweis {lingkaran kelengkungan}{}{}
untuk

\mavergleichskette
{\vergleichskette
{ t
}
{ \in }{ \R
}
{  }{
}
{    }{
}
{   }{
}
}
{}{}{}
dengan syarat $f'(t) \neq 0$, menggunakan rumus-rumus dalam
Lema 3.10.

}
{} {}

\inputaufgabe
{}
{

Tentukan
\definitionsverweis {kelengkungan}{}{}
\definitionsverweis {grafik}{}{}
\maabbdisp {\sin} {\R} {\R
} {}
untuk

\mavergleichskette
{\vergleichskette
{ t
}
{ = }{ { \frac{   \pi }{  2 } }
}
{  }{
}
{    }{
}
{   }{
}
}
{}{}{.}

}
{} {}

\inputaufgabe
{}
{

Tentukan
\definitionsverweis {kelengkungan}{}{}
spiral logaritmik
\maabbeledisp {\gamma} {\R} { \R^2
} { t } { \gamma(t) = e^t \left(  \cos  t   , \,   \sin  t                   \right)
} {,}
untuk setiap

\mavergleichskette
{\vergleichskette
{ t
}
{ \in }{ \R
}
{  }{
}
{    }{
}
{   }{
}
}
{}{}{.}

}
{} {}

Sebuah fungsi
\maabb {f} { I } {\R
} {}
pada interval terbuka

\mavergleichskette
{\vergleichskette
{ I
}
{ \subseteq }{ \R
}
{  }{
}
{    }{
}
{   }{
}
}
{}{}{}
disebut
\definitionswort {analitik}{,}
jika pada setiap titik

\mavergleichskette
{\vergleichskette
{ a
}
{ \in }{ I
}
{  }{
}
{    }{
}
{   }{
}
}
{}{}{}
fungsi tersebut dapat dinyatakan oleh suatu
\definitionsverweis {deret pangkat}{}{.}

Sebuah kurva disebut analitik jika setiap fungsi komponennya analitik.

\inputaufgabe
{}
{

Misalkan
\maabb {\kappa} { I } {\R
} {}
sebuah
\definitionsverweis {fungsi analitik}{}{.}
Buktikan bahwa terdapat kurva analitik
\maabb {\gamma} { I } {\R^2
} {}
yang
\definitionsverweis {kelengkungan}{}{}
di titik $\gamma(t)$ sama dengan $\kappa(t)$.

}
{} {}

\inputaufgabe
{}
{

\bild{ \begin{center}
\includegraphics[width=5.5cm]{ \bildeinlesungsvg {Evolute-parab} {svg} }
\end{center}
\bildtext {} }

\bildlizenz { Evolute-parab.svg } {} {Ag2gaeh} {Commons} {CC-by-sa 4.0} {}

Tentukan
\definitionsverweis {evolut}{}{}
kurva
\maabbeledisp {} {\R} {\R^2
} { t } { \begin{pmatrix}  t  \\ t^2    \end{pmatrix}
} {.}

}
{} {}

\inputaufgabegibtloesung
{}
{

Tentukan
\definitionsverweis {kelengkungan}{}{}
di setiap titik lingkaran satuan yang diberikan secara implisit oleh

\mavergleichskettedisp
{\vergleichskette
{ h(x,y)
}
{ =} { x^2+y^2
}
{ =} { 1
}
{ } {
}
{ } {
}
}
{}{}{}
menggunakan
Lema 3.11
dan medan gradien $h$.

}
{} {}

  \zwischenueberschrift{Soal untuk dikumpulkan}

\inputaufgabe
{2}
{

Tentukan
\definitionsverweis {kelengkungan}{}{}
\definitionsverweis {grafik}{}{}
fungsi eksponensial
\maabbdisp {\exp} {\R} {\R
} {}
untuk setiap

\mavergleichskette
{\vergleichskette
{ t
}
{ \in }{ \R
}
{  }{
}
{    }{
}
{   }{
}
}
{}{}{.}

}
{} {}

\inputaufgabe
{2}
{

Tentukan
\definitionsverweis {kelengkungan}{}{}
spiral Archimedes
\maabbeledisp {\gamma} {\R} { \R^2
} { t } { \gamma(t) = t \left(  \cos  t   , \,   \sin  t                   \right)
} {,}
untuk setiap

\mavergleichskette
{\vergleichskette
{ t
}
{ \in }{ \R
}
{  }{
}
{    }{
}
{   }{
}
}
{}{}{.}

}
{} {}

\inputaufgabe
{4}
{

Berikan sebuah contoh kurva berparametrisasi panjang busur
\maabbdisp {\gamma} {\R} {\R^2
} {}
yang dua kali terdiferensialkan secara kontinu, sedemikian sehingga untuk dua waktu

\mavergleichskette
{\vergleichskette
{ t_0
}
{ \neq }{ t_1
}
{  }{
}
{    }{
}
{   }{
}
}
{}{}{}
berlaku kesamaan

\mavergleichskette
{\vergleichskette
{ \gamma (t_0)
}
{ = }{ \gamma(t_1)
}
{  }{
}
{    }{
}
{   }{
}
}
{}{}{}
dan
\definitionsverweis {lingkaran-lingkaran kelengkungan}{}{}
pada
\mathkor {} {t_0} {dan} {t_1} {}
tidak sama.

}
{} {}

\inputaufgabe
{4}
{

Tentukan
\definitionsverweis {evolut}{}{}
kurva
\maabbeledisp {} {\R} {\R^2
} { t } { \begin{pmatrix}  2 \cos  t  \\ 3 \sin  t    \end{pmatrix}
} {.}
Di titik-titik manakah turunan evolut mempunyai titik nol?

}
{} {}

\inputaufgabe
{4}
{

Tentukan
\definitionsverweis {kelengkungan}{}{}
lingkaran satuan yang diberikan secara implisit oleh

\mavergleichskettedisp
{\vergleichskette
{ h(x,y)
}
{ =} { x^4+y^4
}
{ =} { 1
}
{ } {
}
{ } {
}
}
{}{}{}
menggunakan
Lema 3.11
dan medan gradien $h$ di titik-titik berikut.
\aufzaehlungzwei {
\mavergleichskette
{\vergleichskette
{ P
}
{ = }{ \begin{pmatrix}  1  \\ 0   \end{pmatrix}
}
{  }{
}
{    }{
}
{   }{
}
}
{}{}{,}
} {
\mavergleichskette
{\vergleichskette
{ Q
}
{ = }{ \begin{pmatrix}  { \frac{   1  }{  \sqrt[4]{2}  } }  \\ { \frac{   1  }{  \sqrt[4]{2}  } }     \end{pmatrix}
}
{  }{
}
{    }{
}
{   }{
}
}
{}{}{.}
}

}
{} {}


\section*{Solusi yang disediakan oleh sumber}
\addcontentsline{toc}{section}{Solusi yang disediakan oleh sumber}
\subsection*{Solusi Soal 3.7}
% Authority: German Wikiversity pageid 151168, revid 1095913, 2026-06-15T07:25:24Z.
% Exact UTF-8 wikitext witness SHA-256: c771b7c1ed173ddc05c8c4219c3c34e7e8099fe42079740b292d991b2ce3fb25.
% Expanded German TeX source SHA-256: 7e3437274bf4f79b6a3fa719876b09fdbcee4e10523a215e9e6f4ecb566cdb85.

Nilai-nilainya adalah

\mavergleichskettedisp
{\vergleichskette
{ \gamma(0)
}
{ =} { \begin{pmatrix}  1  \\ 0   \end{pmatrix}
}
{ } {
}
{ } {
}
{ } {
}
}
{}{}{,}

\mavergleichskettedisp
{\vergleichskette
{ \gamma'(0)
}
{ =} { \begin{pmatrix}  0  \\ \omega   \end{pmatrix}
}
{ } {
}
{ } {
}
{ } {
}
}
{}{}{}
dan

\mavergleichskettedisp
{\vergleichskette
{ \gamma^{\prime \prime} (0)
}
{ =} { \begin{pmatrix}  - \omega^2  \\ 0   \end{pmatrix}
}
{ } {
}
{ } {
}
{ } {
}
}
{}{}{.}
Suatu gerak melingkar yang memiliki nilai serta turunan pertama dan kedua yang sama dengan gerak yang diberikan pada

\mavergleichskette
{\vergleichskette
{ t
}
{ = }{ 0
}
{  }{
}
{    }{
}
{   }{
}
}
{}{}{}
harus, khususnya, mempunyai garis singgung yang sama di titik ini. Karena itu, pusat lingkaran harus terletak pada sumbu $x$. Selain itu, pusat tersebut harus berada di sebelah kiri
\mathl{\begin{pmatrix}  1  \\ 0   \end{pmatrix}}{},
karena percepatan mengarah ke kiri. Oleh sebab itu, kita dapat menuliskan gerak melingkar beraturan yang dicari, dengan pusat
\mathl{\begin{pmatrix}  x  \\ 0   \end{pmatrix}}{}
\zusatzklammer {
\mavergleichskettek
{\vergleichskettek
{ x
}
{ < }{  1
}
{  }{
}
{    }{
}
{   }{
}
}
{}{}{}} {} {}
dan jari-jari
\mathl{1-x}{}, sebagai

\mavergleichskettedisp
{\vergleichskette
{  \delta (t)
}
{ =} {  \begin{pmatrix}  x  \\ 0   \end{pmatrix} + (1-x) \begin{pmatrix}  \cos \left(  u t \right)  \\ \sin  \left(  u t \right)     \end{pmatrix}
}
{ } {
}
{ } {
}
{ } {
}
}
{}{}{.}
Kita memperoleh

\mavergleichskettedisp
{\vergleichskette
{ \delta(0)
}
{ =} { \gamma(0)
}
{ } {
}
{ } {
}
{ } {
}
}
{}{}{.}
Selanjutnya,

\mavergleichskettedisp
{\vergleichskette
{ \delta'( 0)
}
{ =} {  (1-x)  \begin{pmatrix}  0  \\ u     \end{pmatrix}
}
{ } {
}
{ } {
}
{ } {
}
}
{}{}{}
dan

\mavergleichskettedisp
{\vergleichskette
{ \delta^{\prime \prime}( 0)
}
{ =} {  (1-x)  \begin{pmatrix}  -u^2  \\ 0     \end{pmatrix}
}
{ } {
}
{ } {
}
{ } {
}
}
{}{}{.}
Jadi, agar syarat-syarat yang diminta terpenuhi, harus berlaku

\mavergleichskettedisp
{\vergleichskette
{ \omega
}
{ =} { (1-x) u
}
{ } {
}
{ } {
}
{ } {
}
}
{}{}{}
dan

\mavergleichskettedisp
{\vergleichskette
{ \omega^2
}
{ =} { (1-x) u^2
}
{ } {
}
{ } {
}
{ } {
}
}
{}{}{.}
Dari sini diperoleh

\mavergleichskettedisp
{\vergleichskette
{ \omega^2
}
{ =} {  \omega u
}
{ } {
}
{ } {
}
{ } {
}
}
{}{}{}
dan dengan demikian

\mavergleichskettedisp
{\vergleichskette
{ u
}
{ =} { \omega
}
{ } {
}
{ } {
}
{ } {
}
}
{}{}{}
serta

\mavergleichskettedisp
{\vergleichskette
{ x
}
{ =} { 0
}
{ } {
}
{ } {
}
{ } {
}
}
{}{}{.}

\subsection*{Solusi Soal 3.16}
\input{generated/worksheet03_exercise16_solution.id.build.tex}

\backmatter
\listoffigures
\chapter*{Atribusi dan Hak Media}
\addcontentsline{toc}{chapter}{Atribusi dan Hak Media}
\section*{Unit 1}
\addcontentsline{toc}{section}{Unit 1}
Empat berkas berikut digunakan dalam Unit 1. Setiap berkas tetap mengikuti lisensinya sendiri; tautan sumber dan lisensi dapat diklik.
\begin{description}
\item[\texttt{3d-function-6.svg}] MartinThoma (karya asli).
\href{https://commons.wikimedia.org/wiki/File:3d-function-6.svg}{Halaman sumber di Wikimedia Commons}; lisensi \href{https://creativecommons.org/licenses/by/3.0}{CC BY 3.0}.
\item[\texttt{Great circle passing through two points.svg}] HaEr48; karya turunan dari Polar angle to spherical side.svg oleh Episcophagus.
\href{https://commons.wikimedia.org/wiki/File:Great_circle_passing_through_two_points.svg}{Halaman sumber di Wikimedia Commons}; lisensi \href{https://creativecommons.org/licenses/by-sa/4.0}{CC BY-SA 4.0}.
\item[\texttt{2019-07-Helix.jpg}] Ag2gaeh (karya asli).
\href{https://commons.wikimedia.org/wiki/File:2019-07-Helix.jpg}{Halaman sumber di Wikimedia Commons}; lisensi \href{https://creativecommons.org/licenses/by-sa/4.0}{CC BY-SA 4.0}.
\item[\texttt{Planned flight map of the Oiseau Blanc.svg}] Pethrus; karya turunan dari BlankMap-World8.svg oleh AMK1211.
\href{https://commons.wikimedia.org/wiki/File:Planned_flight_map_of_the_Oiseau_Blanc.svg}{Halaman sumber di Wikimedia Commons}; lisensi \href{https://creativecommons.org/licenses/by-sa/3.0}{CC BY-SA 3.0}.
\end{description}
Identitas revisi, ukuran, SHA-1 Wikimedia, SHA-256, URL berkas asli, dan hash turunan cetak tercatat dalam manifest media yang disertakan.

\section*{Unit 2}
\addcontentsline{toc}{section}{Unit 2}
Dua berkas berikut digunakan dalam Unit 2. Setiap berkas tetap mengikuti status hak atau lisensinya sendiri; tautan sumber dapat diklik dan tautan lisensi dicantumkan jika tersedia.
\begin{description}
\item[]\texttt{Integral} \texttt{apl} \texttt{rot} \texttt{obsah1.svg}
Pajs; dialihkan dari Wikipedia bahasa Ceko ke Wikimedia Commons.
\href{https://commons.wikimedia.org/wiki/File:Integral_apl_rot_obsah1.svg}{Halaman sumber di Wikimedia Commons}; status hak: domain publik.
\item[]\texttt{Hyperboloid1.png}
Lars H. Rohwedder (RokerHRO; karya asli).
\href{https://commons.wikimedia.org/wiki/File:Hyperboloid1.png}{Halaman sumber di Wikimedia Commons}; status hak: domain publik.
\end{description}
Identitas revisi, ukuran, SHA-1 Wikimedia, SHA-256, URL berkas asli, dan hash turunan cetak tercatat dalam manifest media yang disertakan.

\section*{Unit 3}
\addcontentsline{toc}{section}{Unit 3}
Tiga berkas berikut digunakan dalam Unit 3. Setiap berkas tetap mengikuti status hak atau lisensinya sendiri; tautan sumber dapat diklik dan tautan lisensi dicantumkan jika tersedia.
\begin{description}
\item[]\texttt{Parabola} \texttt{circle.svg}
IkamusumeFan (karya sendiri).
\href{https://commons.wikimedia.org/wiki/File:Parabola_circle.svg}{Halaman sumber di Wikimedia Commons}; lisensi \href{https://creativecommons.org/licenses/by-sa/4.0}{CC BY-SA 4.0}.
\item[]\texttt{Euler} \texttt{spiral.svg}
AdiJapan (karya sendiri).
\href{https://commons.wikimedia.org/wiki/File:Euler_spiral.svg}{Halaman sumber di Wikimedia Commons}; lisensi \href{https://creativecommons.org/licenses/by-sa/3.0}{CC BY-SA 3.0}.
\item[]\texttt{Evolute-parab.svg}
Ag2gaeh (karya sendiri).
\href{https://commons.wikimedia.org/wiki/File:Evolute-parab.svg}{Halaman sumber di Wikimedia Commons}; lisensi \href{https://creativecommons.org/licenses/by-sa/4.0}{CC BY-SA 4.0}.
\end{description}
Identitas revisi, ukuran, SHA-1 Wikimedia, SHA-256, URL berkas asli, dan hash turunan cetak tercatat dalam manifest media yang disertakan.

\chapter*{Lisensi}
\Lizenztext
\end{document}
